\documentclass[11pt]{article}
\usepackage[T1]{fontenc}
\usepackage[utf8]{inputenc}
\usepackage{lmodern,amsmath,amssymb,amsthm,mathtools,microtype,booktabs}
\usepackage[margin=27mm]{geometry}
\usepackage[round,authoryear]{natbib}
\usepackage{xurl}
\usepackage[hidelinks]{hyperref}
\newcommand{\workbenchurl}{https://github.com/KokunoYumeto/arithmetic-geometry-complexity-workbench}
\newcommand{\recordurl}{https://github.com/KokunoYumeto/mathematics-workbenches}
\newcommand{\N}{\mathbb N}
\DeclareMathOperator{\Mon}{Mon}
\DeclareMathOperator{\Arith}{Arith}
\DeclareMathOperator{\Sup}{Sup}
\newtheorem{theorem}{Theorem}[section]
\newtheorem{proposition}[theorem]{Proposition}
\newtheorem{lemma}[theorem]{Lemma}
\newtheorem{corollary}[theorem]{Corollary}
\theoremstyle{definition}
\newtheorem{definition}[theorem]{Definition}
\newtheorem{example}[theorem]{Example}
\title{Monotone circuits for the edge functions of $C_4$-free graphs}
\author{}
\date{25 November 2026}
\begin{document}
\maketitle
\begin{abstract}
Jukna's work on quadratic Boolean functions shows why circuit computation cannot in general be reduced without cost to a single product level. The public P versus NP programme asks whether the edge function of every graph of girth at least five requires exactly the number of edges plus the vertex-cover number minus one binary monotone gates. Starting from the established theorem for at most three missing quadratic witnesses, we prove that four missing induced edges require at least four deleted products, exclude both the optimal six-deletion and optimal seven-deletion cases, and prove the full lower bound whenever four missing witnesses increase the vertex-cover number by four. A full-graph safety argument handles the remaining quadratic circuit with excess one. At cover gain three, we prove the full bound for an optimal quadratic frame when the missing graph is a three-edge path and a disjoint edge and the ambient graph has girth at least five; we also exclude five deletions for a missing triangle and disjoint edge in a four-cycle-free graph. With only one deleted addition, we prove an exponential lower bound on the deletion count in terms of the cover gain. We also give a size-preserving normalization for a specified shared two-level bilinear class. Complete propositional certificates support the finite obstructions. We prove the earlier bounded-weight barrier and one-point constructions in full, and correct the negation budget in the fast convolution construction using Beals, Nishino and Tanaka. The unrestricted conjecture, the remaining four-witness cases and general syntactic bidegree two remain open. The mathematics and text were produced by GPT-6 Astra Pro; separate re-proofs by the same model and exact certificate checks are recorded, and no human has checked the paper.
\end{abstract}

\section{The question and the model}
The purpose of this paper is to isolate the first unresolved witness configurations in a concrete monotone-circuit lower-bound problem.
Let $G=(V,E)$ be a finite simple graph with $E\ne\varnothing$. Its edge function is
\[
 f_G(x)=\bigvee_{\{u,v\}\in E}(x_u\wedge x_v).
\]
Let $\tau(G)$ denote the minimum size of a vertex cover. A circuit is a finite acyclic directed graph of binary AND and OR gates, with variables at its inputs, one designated output, and unrestricted reuse of wires. Its size is the number of gates; inputs have cost zero. Constants are not part of the model. Write $\Mon(f_G)$ for its minimum size.

The conjecture considered in the public P versus NP programme is
\begin{equation}\label{eq:conjecture}
 \Mon(f_G)=|E|+\tau(G)-1
\end{equation}
when $G$ has girth at least five. The proposed lower bound is also studied for all graphs without a four-cycle, including triangles. These are different hypotheses, and every theorem below states which it uses.

\begin{proposition}\label{prop:upper}
For every finite simple graph with at least one edge,
$\Mon(f_G)\leq|E|+\tau(G)-1$.
\end{proposition}
\begin{proof}
Choose a minimum cover $K$ and assign each edge to one of its endpoints in $K$. Every vertex of $K$ has an edge to a vertex outside $K$, since otherwise it could be removed from the cover; assign that edge to it. Thus every assigned star is nonempty. If its leaf set has size $d_v$, compute
$x_v\wedge\bigvee_{u\text{ assigned to }v}x_u$ using $d_v-1$ OR gates and one AND gate. Finally combine the $|K|$ star outputs with $|K|-1$ OR gates. As $\sum_vd_v=|E|$, the total is the claimed bound, and the computed function is exactly $f_G$.
\end{proof}

The construction matters because an improvement must exploit Boolean identities absent from ordinary nonnegative polynomial computation. A square $x_u^2$ and a variable $x_u$ agree on Boolean inputs, while their formal degrees differ. The lower-bound problem must charge these identifications together with wire sharing.

Jukna's counterexample to the general single-level conjecture \citep{JuknaSingle} prevents a cost-preserving flattening argument for arbitrary quadratic Boolean functions. Formula lower bounds also require care: a formula is a tree, whereas the circuits here can reuse a computation. Neither comparison settles \eqref{eq:conjecture}.

\subsection{Main theorems and their meaning}
For a positive arithmetic polynomial $P$, say $P\in\Sup(G)$ if every monomial support contains an edge of $G$ and each edge is the exact support of some monomial. Replacing the OR and AND labels of a variable-input circuit by addition and multiplication gives such a polynomial exactly when the circuit computes $f_G$, by Jukna's support criterion \citep[Lemma~1, p.~7]{JuknaRead}.
Let $H$ be the graph of edges whose squarefree quadratic monomials occur in $P$. The minimum-degree construction of Section~\ref{sec:four} produces its quadratic circuit $C_2$ without adding gates. Put
\[
 d=\tau(G)-\tau(H),\quad D=|C|-|C_2|,\quad
 e=|C_2|-\bigl(|E(H)|+\tau(H)-1\bigr).
\]
A deleted product or addition means a gate absent from this construction after unused gates have been removed.

\begin{theorem}[The four-witness reductions]\label{thm:four-summary}
Let $G$ be a finite simple graph with no four-cycle and let a positive variable-input binary circuit $C$ compute $P\in\Sup(G)$ with nonempty quadratic support $H$.
If eight vertices induce four disjoint edges in $G$ and no edge in $H$, at least four product gates are deleted in forming $C_2$.
If $|E(G)\setminus E(H)|=4$ and $|C|<|E(G)|+\tau(G)-1$, the only remaining numerical possibilities are
\[
\begin{array}{c|l}
 d&(D,e)\\\hline
 2&(5,0)\\
 3&(5,0),(5,1),(6,0).
\end{array}
\]
In particular, every circuit with four missing witnesses and cover gain four satisfies the full star lower bound. More generally, if the four missing edges form an induced matching and six gates are deleted, four of them products, then $\tau(G)\le K+2$, where $K$ is the number of products in $C_2$; this again implies the full star lower bound. In the row $(d,D,e)=(3,5,0)$, the missing-edge graph cannot be $P_4\sqcup K_2$ or $K_3\sqcup K_2$, where $P_4$ denotes the path with three edges. If $G$ has girth at least five, the shape $P_4\sqcup K_2$ is also excluded in the row $(d,D,e)=(3,6,0)$; hence that shape satisfies the full bound at cover gain three whenever its quadratic frame is optimal.
\end{theorem}
The inherited numerical alternatives are \citet[Corollary~5.34]{Companion}. The additional exclusions and the four-product obstruction are proved in Section~\ref{sec:four}. Thus a four-witness counterexample must satisfy the displayed strict restrictions; the theorem does not exclude every such counterexample. This narrows the first case beyond the established at-most-three theorem without asserting the full circuit conjecture.

The restriction with one deleted addition does not require four missing witnesses or an optimal quadratic frame.
\begin{theorem}[One deleted addition]\label{thm:pure-summary}
Let $G$ be a finite simple graph without a four-cycle, and let a positive variable-input binary circuit computing $P\in\Sup(G)$ have nonempty quadratic support $H$. If exactly one deleted gate is an addition, then $G[V(G)\setminus Z]$ is a matching with possible isolated vertices for every vertex cover $Z$ of $H$, and $D\ge2^d-1$. If $G$ has girth at least five, the missing-edge graph itself is a matching.
\end{theorem}
The proof is Theorem~\ref{thm:purecut}. In particular, all remaining cover-gain-three alternatives require at least two deleted additions.

\begin{theorem}[A class at two product levels]\label{thm:bilinear-summary}
Let $G=(X\sqcup Y,E)$ be a nonempty bipartite graph without a four-cycle. Suppose a binary monotone circuit computing $f_G$ has first products $R_i=A_i(X)B_i(Y)$, where $A_i,B_i$ are pure-side linear disjunctions, second products multiplying disjunctions of the $R_i$, and an output disjunction of first and second products. All disjunction networks and operands may be shared. Call $R_i$ unsafe if $A_i\times B_i$ contains a nonedge. If every unsafe rectangle has at most one row $x\in A_i$ satisfying $B_i\subseteq N(x)$, the circuit has at least $|E|+\tau(G)-1$ gates. The conclusion also holds under the symmetric column condition.
\end{theorem}
The proof in Section~\ref{sec:degree-two} changes the shared circuit into a quadratic circuit without growth and then applies the established read-one theorem. It covers every such architecture whose unsafe rectangles all have at least two columns. General bidegree $(2,2)$ also allows mixed-degree operands and more shared intermediate expressions, and is not settled by this statement.

The other results explain two limits of proposed approaches. Theorem~\ref{thm:weightbarrier} constructs small circuits correct on all edges and on every sufficiently small independent set, so a large lower bound must use larger negative inputs. In contrast, the convolution construction computes all projective-plane line outputs with nearly linear general-circuit size, and the checked negation-reduction theorem places this upper bound at a specified logarithmic NOT budget. These are statements about different input promises and different gate sets.

\section{Proof overview}
The lower-bound argument separates the quadratic polynomial carried to the output from the products and additions deleted by minimum-degree extraction. Earlier cover-exchange arguments reduce four missing witnesses to a short list of deletion counts. An induced matching then turns Boolean correctness into twenty decisive tests: four positive edges and sixteen maximal independent transversals. We describe a larger Boolean template than any circuit in the relevant configuration and exclude it by a complete reverse-unit-propagation certificate. For seven deletions, clause counting and a confinement argument remove the five-product case. Four products leave one, two or three visible cuts. A complete enumeration allowing all wire sharing gives at most thirteen CNF clauses for the one-cut branch, below the sixteen needed for a four-edge matching. The two- and three-cut branches reduce to separately verified finite obstructions. Together these arguments exclude the entire optimal seven-deletion case. For a quadratic frame with excess one, the final functional possibility has two quartic contributions. Full-graph safety and four-cycle-freeness force their entire insertion contexts to be disjoint pairs of missing endpoints. Those four vertices replace at least two quadratic star centers, giving a cover too small for cover gain four.

With one deleted addition, the high branch is a conjunction of old linear clauses. A two-edge path outside an old cover would force a minimal-cover clause containing its two outer vertices, while four-cycle-freeness excludes that pair from every old clause and from the single context. The residual graph must therefore be a matching; counting its minimal covers gives the exponential deletion bound.

For a missing three-edge path and a disjoint edge, all six endpoints lie outside every minimum old cover when the cover gain is three. An optimal quadratic frame then prevents any nonzero direct context from lying in a minimum cover of the missing graph. This restriction, together with the prohibition on old clauses containing the outer vertices of a missing two-edge path, excludes five deletions. For six deletions in girth at least five, distinct visible cuts satisfy a further common-neighbor restriction on the singleton assignments accepted by their high values. Seven explicit finite classes cover all remaining product and addition profiles and exclude them.

At bidegree two the main difficulty is sharing: expanding products into rectangles can multiply the number of gates. The simultaneous replacement theorem keeps the existing linear wires, replaces each unsafe rectangle by its unique possible safe row, and changes the second products to disjunctions. Safety and preservation of every edge witness prove correctness without duplicating wires.

The bounded-weight construction uses coordinate changes preserving affine incidence and takes the conjunction of several coordinate tests. Its probability estimate concerns independent sets, so it does not prove an exact circuit for the graph. The negation results use a different mechanism: cyclic convolution computes all Singer line outputs, while a general simulation theorem controls the NOT budget. Neither mechanism provides the missing monotone lower bound.

\section{The polynomial witness condition}
This section makes the distinction between a quadratic witness and a minimal Boolean support explicit.
Replace every OR by addition and every AND by multiplication in a variable-input circuit $C$. The resulting polynomial $P_C\in\N[x_v:v\in V]$ has positive coefficients on all its occurring monomials. Define $\Sup(G)$ to be the class of such nonzero polynomials for which every monomial support contains an edge of $G$, and every edge is the exact support of some monomial.

\begin{proposition}\label{prop:support}
The Boolean circuit $C$ computes $f_G$ if and only if $P_C\in\Sup(G)$. Consequently
\[
 \Mon(f_G)=\min_{P\in\Sup(G)}\Arith(P),
\]
where $\Arith(P)$ is the minimum number of binary addition and multiplication gates in a variable-input circuit computing $P$ over $\N$.
\end{proposition}
This is the specialization of Jukna's support criterion to quadratic minimal true assignments \citep[Lemma~1, p.~7 of the arXiv manuscript]{JuknaRead}. The equality follows because changing the gate labels between the Boolean and nonnegative arithmetic interpretations preserves the circuit graph. We use this established criterion throughout.

For $P\in\Sup(G)$, let $E_2(P)$ be the edges $uv$ whose quadratic monomial $x_ux_v$ occurs. An edge outside $E_2(P)$ has a witness $x_u^ax_v^b$ with $a,b\geq1$ and $a+b>2$. In the terminology of Jukna's witness model, a read-1 circuit for this quadratic function has $E_2(P_C)=E$. This does not mean that a variable has only one occurrence in the circuit.
The one-edge polynomial $x_u^2x_v$ is the smallest example showing why an arbitrary witness cannot be recovered from its quadratic part.

\section{Established exact cases and the remaining obstruction}
This section specifies the lower bounds used as inputs to the four-witness analysis.
The public proof graph's Lean theorem L-meta-2237 proves the read-1 lower bound for every nonempty finite simple graph without a four-cycle \citep{Fleet}.

The corresponding formal statement is \mbox{\nolinkurl{PEqualNP.Outside.L_meta_2237}}.
Combining it with the star circuit gives
\[
 \min_{c_e\in\mathbb N_{>0}}\Arith\!\left(\sum_{e=uv\in E}c_e x_ux_v\right)
 =\Arith\!\left(\sum_{e=uv\in E}x_ux_v\right)
 =\operatorname{Read}_1(f_G)=|E|+\tau(G)-1.
\]
Here read-1 means that every minimal true assignment has a squarefree witness, as in \citet[Definition~1, p.~8; Theorem~1]{JuknaRead}. It is distinct from a bound on wire fan-out.

We also use the following two results from the companion paper \citep[Theorem~1.3]{Companion} and \citep[Theorem~1.1]{OneSidedCompanion}. They are mathematical theorems from those companion manuscripts; they are not asserted here to be Lean theorems.
\begin{theorem}[Established witness and one-sided bounds]\label{thm:granted}
Let $G$ be a finite simple graph with $E\ne\varnothing$ and no four-cycle.
If a variable-input positive arithmetic circuit $C$ computes $P\in\Sup(G)$ and at most three edges lack their quadratic monomial in $P$, then
$|C|\geq|E|+\tau(G)-1$.
If $G=(X\sqcup Y,E)$ is bipartite and every monomial of $P$ contains at most one distinct variable from $X$, the same bound holds. In the latter case, the separate counts are at least $|E|-1$ additions and at least $\tau(G)$ products. The same assertion holds with $X$ and $Y$ exchanged.
Both total bounds are attained by a star circuit.
\end{theorem}
These are inputs rather than new proofs. In particular, an improving circuit must miss at least four quadratic witnesses, and its output must contain a monomial mixing distinct variables from $X$ and a monomial mixing distinct variables from $Y$; the two monomials need not coincide.
The syntactic degree-at-most-two condition on each side restricts a larger class than the one-sided theorem and remains a separate question.

\section{Four missing witnesses and a smaller list of cases}\label{sec:four}
This section proves the full lower bound when four missing witnesses increase the cover number by four. The argument begins with a product-count obstruction, then treats the possible ways the higher-degree contributions enter the quadratic circuit. Let $H$ be the quadratic support graph, and let $C_2$ be the circuit obtained by minimum-degree extraction followed by pruning. Define
\[
 d=\tau(G)-\tau(H),\qquad D=|C|-|C_2|,\qquad
 e=|C_2|-(|E(H)|+\tau(H)-1).
\]
Thus $d$ is the increase in cover number, $D$ counts deleted gates, and $e$ is the excess size of the quadratic circuit. In the four-witness case, the existing reduction leaves precisely the rows of \citet[Corollary~5.34]{Companion}. We remove the pair $(D,e)=(6,0)$ from its $d=4$ row.

Here is the extraction used throughout the argument. The minimum degree
$\delta$ is one at a variable, is the sum of the input degrees at a
product, and is their minimum at an addition. Starting from the output,
select both inputs at a product or an equal-degree addition, and only the
lower-degree input at an unequal-degree addition. Bypass each selected
unequal-degree addition and discard unused gates. The resulting circuit
$C_2$, called the \emph{quadratic frame}, computes the degree-two part of
$P$; positivity prevents cancellation. A \emph{private} gate is absent
from this retained circuit. An unequal-degree addition is a
\emph{cut}; its lower- and higher-degree operands are its low and high
inputs. A cut is \emph{visible} if its own node belongs to the selected
output subgraph before bypassing. Reuse of its low input elsewhere does
not make the cut visible.

For the Boolean arguments, a function is \emph{safe} for $G$ if every
assignment on which it is true contains an edge of $G$. It
\emph{witnesses} an edge if it is true on the assignment consisting of
that edge's two endpoints. A positive clause is a disjunction of input
variables; the empty clause is zero. An affine positive clause may also
have a constant term, in which case a positive constant makes it one.
A conjunctive normal form (CNF) is a conjunction of such clauses.
These constants are permitted in the relaxed Boolean templates and in
restricted circuits, even though the original circuit has only variable
inputs. Positive arithmetic linear forms give these clauses on Boolean
supports.


The following finite lemma is proved by an explicit propositional certificate. It is used only as an obstruction on a larger class: the template deliberately omits the exact-frame, optimality and pruning restrictions that a circuit would satisfy. Thus it is not a counterexample search, and its unsatisfiability cannot discard a legitimate circuit.

\begin{lemma}[A three-product Boolean template]\label{lem:threeprivatefinite}
For $i=0,1,2$, let
\[
 g_i=A_i(x,g_0,\ldots,g_{i-1})\wedge B_i(x,g_0,\ldots,g_{i-1}),
\]
where $x=(x_0,\ldots,x_7)$ and $A_i,B_i$ are arbitrary disjunctions of subsets of the displayed arguments; the empty disjunction is allowed.
Let $L_i$ be arbitrary affine positive clauses in $x$, including zero and one, and let $c_{ij}\in\{0,1\}$.
Then
\[
 \bigvee_i(g_i\wedge L_i)\ \vee\!
 \bigvee_{i<j:c_{ij}=1}(g_i\wedge g_j)
 \ne\bigvee_{a=0}^3(x_{2a}\wedge x_{2a+1}).
\]
\end{lemma}
\begin{proof}[Proof by finite certificate]
Use one selector for each possible argument in each input clause, each possible variable or constant in each $L_i$, and each cross term. There are $54+27+3=84$ selectors. For fixed selectors the function is monotone, so equality with the matching is equivalent to truth on its four edge assignments and falsity on its sixteen maximal independent endpoint transversals.
The program \emph{Private-product obstruction} expresses evaluations on these twenty assignments by exact Tseitin equivalences: a selected argument is the conjunction of its selector and its value, each input is their disjunction, and each product and output gate is encoded in both directions.
The resulting CNF has $2{,}085$ variables and $7{,}421$ clauses.

Its DRUP proof contains $43{,}308$ additions and $41{,}308$ deletions. An independent checker using only the C++ standard library verifies every addition by reverse unit propagation and derives the empty clause. Deletions are ignored, soundly retaining only original or already implied clauses. The check took $15.18$ seconds. The encoder, proof and standard-library checker are available through the workbench. Reverse unit propagation means that adjoining the negation of the proposed clause and propagating forced literals yields a contradiction. Therefore each accepted clause follows from the earlier clauses; induction ends in a proof of inconsistency. The complete check is a propositional certificate.
\end{proof}

\begin{proposition}[Four missing induced edges require four deleted products]\label{prop:fourprivateproducts}
Let a positive variable-input binary circuit compute $P\in\Sup(G)$ with nonempty quadratic support $H$. Suppose eight vertices induce precisely four disjoint edges of $G$ and no edge of $H$. Minimum-degree extraction and pruning delete at least four product gates. There is no bound on deleted additions and no quadratic-optimality or cover-number hypothesis.
\end{proposition}
\begin{proof}
Suppose at most three products are deleted. Apply the full-carrier erasure from \citet[proof of Lemma~5.26, first three paragraphs]{Companion}: protect the original minimum-degree-two nodes with addition-only output paths, and erase their off-carrier uses while their \emph{whole current values} still reach the output additively. The erased multiples are Boolean-absorbed. This preserves the exact extracted quadratic circuit, the Boolean output and the bound of three deleted products. It does not replace a carrier by its quadratic part.

No carrier now feeds its complement. Every product in that complement is deleted, because a quadratic product that survives extraction has an addition-only output path and is itself protected. Order those private products as $g_0,g_1,g_2$, padding with unused zero expressions if necessary. Induction through the private addition gates shows that each input to product $i$ is a nonnegative linear combination of original variables and earlier $g_j$. Booleanly this is exactly the arbitrary input clause permitted by Lemma~\ref{lem:threeprivatefinite}; arbitrary sharing is included.

The operands of a quadratic-frame product have the form
\[
 A+\sum_i a_i g_i,\qquad B+\sum_i b_i g_i,
\]
with $A,B$ linear forms and all coefficients nonnegative. Expand these quadratic-frame products, leaving the private products unexpanded. The protected output network uses additions alone, so summing gives
\[
 P_2+\sum_i g_iL_i+\sum_{i\leq j}c_{ij}g_ig_j,
\]
where $L_i$ are affine nonnegative linear forms in the original variables. Direct private-product contributions give their constant terms. No input variable can reach the output additively because the output has minimum degree two. The input-only contribution is precisely $P_2$: the quadratic-frame products multiply their original degree-one parts, and each private product has minimum degree at least two.

Set every variable outside the specified eight vertices to zero. Then $P_2$ vanishes, and the Boolean output is exactly their four-edge matching. Boolean idempotence absorbs $g_i^2$ into the constant part of $L_i$, and positive coefficients become presence bits. The result is the forbidden template of Lemma~\ref{lem:threeprivatefinite}, a contradiction.
\end{proof}

\begin{corollary}[The optimal six-deletion case]
Among the four-missing-witness alternatives, $(d,D,e)=(4,6,0)$ is impossible. The remaining necessary pairs are
\[
\begin{array}{c|l}
d&(D,e)\\\hline
2&(5,0)\\
3&(5,0),(5,1),(6,0)\\
4&(6,1),(7,0).
\end{array}
\]
\end{corollary}
\begin{proof}
For $d=4$, the missing edges form an induced matching with no old edge on their endpoints. Full-carrier erasure preserves the optimal quadratic frame, and \citet[Lemma~5.33]{Companion} forces all six deleted gates to remain. Proposition~\ref{prop:fourprivateproducts} gives at least four private products. At least one deleted addition is an unequal-degree cut, since otherwise the useful circuit is homogeneous and its degree-two output has no missing witness. Thus the product/addition profile is $(4,2)$ or $(5,1)$.

The four-product case is excluded by \citet[Lemma~5.35]{Companion}. In the five-product case there is a unique cut. Its high input is a connected multiplication DAG with at most five products and hence at most six distinct old linear boundary inputs: ten operand occurrences include at least four internal uses. Its scalar or linear direct quadratic-frame context adds at most one factor. After restriction and Boolean idempotence the matching would therefore be a conjunction of at most seven affine clauses.
Such a representation requires at least sixteen clauses: every nonconstant clause meets each of the four edge assignments and rejects at most one independent endpoint transversal, whereas all sixteen transversals must be rejected. This contradiction completes the proof.
\end{proof}

The same argument leaves only four products and two additions when $(d,D,e)=(4,6,1)$, and only the product/addition profiles $(4,3)$ and $(5,2)$ when $(d,D,e)=(4,7,0)$. These restrictions do not assert the existence of any remaining circuit. The full four-missing-witness theorem remains open here.

\subsection{One deleted addition and exponential cover gain}
A single deleted addition imposes a stronger restriction than the small-profile clause counts alone suggest.

\begin{theorem}[A single deleted addition]\label{thm:purecut}
Let $G$ be a finite simple graph without a four-cycle, and let a positive variable-input binary circuit compute $P\in\Sup(G)$ with nonempty quadratic support $H$. Suppose exactly one of the gates deleted by minimum-degree extraction and pruning is an addition. Then, for every vertex cover $Z$ of $H$, the induced graph $G[V(G)\setminus Z]$ is a matching with possible isolated vertices. Moreover,
\[
 D\ge 2^d-1,\qquad d=\tau(G)-\tau(H).
\]
If $G$ has girth at least five, the missing-edge graph $E(G)\setminus E(H)$ itself is a matching.
\end{theorem}
\begin{proof}
If no edge is missing, the induced graph outside any $H$-cover is empty and the conclusions are immediate. Otherwise apply full-carrier erasure. The sole deleted addition cannot disappear, since its disappearance would leave no route for a missing witness to enter the selected quadratic output. It is the unique visible cut, and every useful private product is ancestral to its high operand $g$. Old quadratic carriers are unavailable as off-carrier inputs. Thus $g$ is a multiplication DAG over old linear boundary values. Booleanly it is a conjunction of their positive clauses.

The independent-symbol quadratic frame is affine in this one high signal. If a degree-one symbol occurred squared, substituting its nonzero old linear low part would give a square in the quadratic output. A degree-two symbol has only scalar context. Hence, after setting any $H$-cover $Z$ to zero, the quadratic part disappears and the output is $gL$, where $L$ is zero, one, or a positive linear clause. This is a positive CNF for the edge function of $G[T]$, where $T=V(G)\setminus Z$.

Suppose $u,w,v\in T$ form a two-edge path, with $u\ne v$. No old boundary clause can contain both $u$ and $v$. A clause containing both variables uses a retained linear-addition gate. That gate reaches a quadratic product, so some opposite-factor variable $h$ is an $H$-neighbor of both. Because $Z$ covers $H$ and $u,v\in T$, the vertex $h$ lies in $Z$ and is distinct from $u,w,v$. The four edges $uw,wv,vh,hu$ would form a four-cycle. The same argument applies to a nonconstant direct context $L$: for its nonzero old low part $A$, the positive product $AL$ occurs in the quadratic frame, so a variable in $A$ would give the same common old neighbor.

Extend $\{w\}$ to a maximal independent set $I$ of $G[T]$. Its complement $U=T\setminus I$ is an inclusion-minimal vertex cover containing $u,v$. Since the CNF $gL$ is false on $I$, some factor clause has support contained in $U$. Every clause of a positive CNF representing an edge function is itself a vertex cover, because it must be true on every edge assignment. Minimality of $U$ forces that clause's support to equal $U$. It would contain $u,v$, although no old factor or the single context can do so. This contradiction proves that $G[T]$ has maximum degree at most one.

Choose $Z$ minimum and let $m$ be the number of edges of the resulting matching. Adding one endpoint of each to $Z$ covers $G$, so $m\ge d$. If $r$ private products remain after erasure, their connected single-output multiplication DAG has at most $r+1$ distinct old boundary sources. Indeed, its $r$ binary nodes have $2r$ input occurrences, at least $r-1$ of which join product nodes; at most $r+1$ boundary occurrences remain. Reuse of a source reduces the number of distinct factors. The context adds at most one clause, so $gL$ has at most $r+2$ clauses.

A matching of $m$ edges has $2^m$ minimal vertex covers. Each is the complement of a maximal independent assignment and must occur as a clause in any positive CNF representing the matching, by the preceding minimality argument. Therefore
\[
 2^d\le2^m\le r+2\le D+1,
\]
since erasure cannot increase the original $D-1$ deleted products. This proves the displayed lower bound without any quadratic-optimality hypothesis.

Finally suppose $G$ is also triangle-free. If two missing edges $uw,wv$ met, there would be no edge $uv$. The three vertices would contain no $H$-edge, so $V(G)\setminus\{u,w,v\}$ would be an $H$-cover. The first conclusion would contradict the two-edge path on its complement. Thus the missing edges are pairwise disjoint.
\end{proof}

In particular, cover gain three with five or six deletions requires at least two deleted additions, whatever the quadratic excess and the shape of the missing edges. This applies to every remaining cover-gain-three alternative in Theorem~\ref{thm:four-summary}.

\subsection{Five products and two additions}
This subsection excludes the five-product profile in the optimal seven-deletion case. Two support lemmas control its analytic branches before the remaining branches are certified.

\begin{lemma}[Shared-factor confinement]\label{lem:sharedfiveconfinement}
Suppose the restricted graph is a matching, possibly with isolated vertices. Let $g$ be a product of $b\leq5$ positive affine clauses. If $gL$ is safe and witnesses two matching edges, with $L$ one more affine clause, then $g$ is false on every other matching edge assignment.
\end{lemma}
\begin{proof}
If $g$ were true on a third edge $\{u,v\}$, one endpoint $a$ would make at least $\lceil b/2\rceil$ of its clauses true; a constant-one clause is true on both endpoints. Adjoin $a$ to each of the four independent transversals of the two witnessed edges. Every remaining clause meets both witnessed edges and rejects at most one of these transversals. Yet all four extended assignments are independent, and only $b+1-\lceil b/2\rceil\leq3$ clauses remain unsatisfied by $a$. This is impossible.
\end{proof}

The next support fact is the cubic rigidity used below. It is the second
matching fact in the proof of \citet[Lemma~5.26]{Companion}; we include
its proof so the confinement argument can be followed here.
\begin{lemma}[Cubic rigidity]\label{lem:cubicrigidity}
Let $ab$ and $xy$ be distinct edges of a matching, with any number of
isolated vertices allowed. Suppose the product of affine positive
clauses $H_1H_2K$ is safe and witnesses $ab$, while both $H_1$ and $H_2$
are true on the edge assignment $xy$. After interchanging endpoints and
the first two factors, its clauses are exactly
\[
 K=a,\qquad H_1=b\vee x,\qquad H_2=b\vee y.
\]
\end{lemma}
\begin{proof}
A safe product witnessing $r$ matching edges needs at least $2^r$
clauses: every clause meets each edge and can reject at most one of
their independent endpoint transversals. Consequently the cubic cannot
witness $xy$, and a two-clause product cannot witness both edges.
If $K=1$, the latter observation gives a contradiction. If $H_1=1$,
safety and the witness $ab$ force the other two clauses to be the
opposite singleton endpoints of $ab$, so $H_2$ cannot be true on $xy$;
the case $H_2=1$ is symmetric. No clause is zero because $ab$ is a
witness. Thus all clauses are nonconstant.

The clause $K$ misses $x,y$. It cannot contain both $a,b$: on the four
independent transversals of these two edges it would always be true,
while $H_1,H_2$ could reject at most two of them. Rename endpoints so
$K\cap\{a,b\}=\{a\}$. On the two transversals $\{a,x\}$ and
$\{a,y\}$, the first two clauses must both miss $a$, contain $b$, and
choose opposite singleton endpoints of $xy$. They therefore restrict
to $b\vee x$ and $b\vee y$. An outside variable $z$ in $K$ would make
the independent assignment $\{b,z\}$ true. One in $H_1$ would make
$\{a,y,z\}$ true, and one in $H_2$ would make $\{a,x,z\}$ true.
Safety excludes every such variable, proving the exact supports.
\end{proof}

\begin{proposition}[The five-product seven-deletion profile]\label{prop:fiveproductseven}
Suppose exactly four quadratic witnesses are missing, $d=4$, and the extracted quadratic circuit is optimal. Seven deleted gates cannot consist of five products and two additions. Consequently the remaining optimal seven-deletion case has exactly four deleted products and three deleted additions.
\end{proposition}
\begin{proof}
First apply full-carrier erasure to the original circuit, before any variable restriction. No private gate can disappear, since that would leave at most six deletions in the already excluded optimal-frame case. Fix the assigned minimum cover $Z$ and retain all isolated vertices when restricting to $T=V\setminus Z$. The restricted graph is the four-edge matching together with those isolated vertices. As in the earlier cover exchange, a nonzero circuit's direct linear context cannot be supported on a transversal of the matching without an isolated vertex: replacing its perturbed center by its outside neighbors, and covering the remaining matching edges, would increase the cover by at most three. Zero and scalar contexts remain allowed.

The gate counts below refer to that original erased circuit; gates that become zero or redundant after restriction remain in the count. If only one cut is visible, its high-input computation has five products and at most one private addition, with at most seven distinct old linear boundary inputs. After idempotence and extraction of common factors its Boolean value is $C_0(A_0\vee B_0)$, where these are products of altogether at most seven clauses. Its scalar context gives at most twelve CNF clauses, and a linear context gives at most thirteen: the clause count is $ab+c$, with total factor count at most seven, or at most eight with $c\geq1$. Both are below sixteen. This also covers a hidden cut or a private equal-degree addition.

Thus both additions are visible cuts. Their high inputs are private-product outputs $g,f$, and the independent-symbol quadratic-frame identity is
\[
 F=gL\vee fM\vee cgf,\qquad c\in\{0,1\}.
\]
The circuit contexts $L,M$ have the restriction stated above. Let $k$ be the number of private products in the original erased circuit ancestral to the first high signal $g$. Every other private product feeds $f$.

If $k=5$, then $f$ is $g$ or one of its ancestors. Write $g=fA$. Their distinct old factors number at most six, and $F=f(AL\vee M\vee cA)$ has at most seven CNF clauses after distribution and absorption.

Suppose $k=4$. Then $g$ has at most five old factors and one product remains for $f$. If $f$ directly reads a proper private ancestor $u$ of $g$, write $g=uW$, with $r\geq1$ factors in $u$, $t$ in $W$, and $r+t\leq5$. If both inputs of $f$ are private ancestors, $f$ divides $g$ and the ancestor argument applies. Otherwise $f=uA$ or $uA\vee g$. Factoring $u$ leaves
\[
 WL\vee AM\vee cWA,
 \qquad W(L\vee M\vee c)\vee AM.
\]
Their CNFs use at most $2(t+1)$ clauses, hence at most $r+2(t+1)\leq11$ including $u$.

In every other case $f=AB$, $AB\vee gB$, or $AB\vee g$, including zero old parts. A safe product with at most seven clauses has at most two matching witnesses, and a safe cubic has at most one. For $f=AB,c=0$, this immediately bounds the union by three. For $c=1$, put $P=gL$, $Q=ABM$, and $R=gAB$. If $R$ has no edge witness, the union has at most three. Otherwise Lemma~\ref{lem:cubicrigidity} implies that a $Q$-witness not witnessed by $R$ would force the direct context $M$ to be a singleton, which is forbidden. Thus $Q$ adds no witness outside $R$. If $P$ has at most one witness the union has at most three; if it has two, Lemma~\ref{lem:sharedfiveconfinement} confines all $R$-witnesses to those two.

For $f=AB\vee gB,c=1$, absorption gives $g(L\vee B)\vee ABM$, with at most three witnesses. For $c=0$, use $P=gL$, $Q=ABM$, $R=gBM$. An outside $Q$-witness would make $B,M$ true on an $R$-edge; Lemma~\ref{lem:cubicrigidity} then makes the direct context $M$ a two-element transversal of those edges. This is forbidden. The same confinement and counting argument applies. Finally, for $f=AB\vee g$ the output is $g(L\vee M)\vee ABM$ when $c=0$, and $g\vee ABM$ when $c=1$; both have at most three witnesses. This completes $k=4$.

For $k=1,2,3$ we use the included finite certificates. Topologically order the $k$ ancestors of $g$ first, ending in $g$, and the remaining products ending in $f$. Before the first cut, each product input is an old linear clause or one earlier private-product output. After that cut it is one earlier private-product output or a clause $A\vee ag$, where $A$ is old linear and $a\in\{0,1\}$. Adding a proper private ancestor to an old clause would require another deleted addition; both available additions are already visible cuts. The full-carrier erasure excludes quadratic-carrier inputs.

Identify all isolated variables with one extra variable $z$. The target remains the matching, and every allowed nonzero nonconstant context contains $z$ or both endpoints of some matching edge. The encoder permits all these contexts and all the input modes just described, without imposing further frame or pruning restrictions. It imposes truth on the four edges with $z=0$ and falsity on the sixteen maximal independent transversals with $z=1$, which by monotonicity is equivalent to the full target.

The three CNFs have respectively $(3700,12897)$, $(3658,12737)$ and $(3616,12577)$ variables and clauses. Complete DRUP proofs are included. The independently compiled DRAT-trim checker verified all three in RUP-only mode, with zero RAT steps, in $87.10$, $104.36$ and $163.31$ seconds. Their hashes, complete compressed proof parts, checker logs, original encoder and reproduction instructions accompany the paper. These certificates exclude $k=1,2,3$ and complete the proof.
\end{proof}

It remains to exclude four private products and three additions in the optimal seven-deletion case. The following subsections do so by the number of visible cuts. The one-excess six-deletion frame and the lower-cover-gain alternatives require separate arguments.

\subsection{Three visible cuts in the seven-deletion case}
This subsection excludes the branch in which all three deleted additions enter the quadratic frame directly, so all three cuts are visible in the sense defined at the start of Section~\ref{sec:four}.

\begin{proposition}[Three visible cuts]\label{prop:threevisible}
Suppose exactly four quadratic witnesses are missing, $d=4$, the extracted quadratic circuit is optimal, and seven deleted gates consist of four products and three additions. The three additions cannot all be visible cuts.
\end{proposition}
\begin{proof}
Apply full-carrier erasure before restricting variables. Every private gate must remain: an erasure or deletion of an unused private gate would give at most six deletions without changing the extracted frame, contrary to the six-deletion exclusion. Choose the assigned minimum cover $Z$ of $H$. After setting its variables to zero, the graph on $T=V(G)\setminus Z$ is the four-edge matching with isolated vertices. Original private gates that become zero or redundant on this restriction stay in the count.

Each cut's high operand is a private-product output. To see this, first note that an old degree-one value is too small to be a high operand. A visible cut of minimum degree one is likewise too small. A selected cut of minimum degree two has an addition-only path to the quadratic output, since any product on that path would raise the minimum degree. Its full value is therefore a protected quadratic carrier. Full-carrier erasure removes every off-carrier use of such a value; it cannot supply the high operand of a degree-one cut. It cannot supply the high operand of a degree-two cut either, because that operand has degree strictly greater than two. The same reasoning excludes old quadratic values. All three deleted additions are already the visible cuts, so no other kind of high operand remains.

Write $p_1,\ldots,p_4$ for the private products in topological order, and mark those selected as high operands. Every operand of a private product is either an earlier private-product output or an old positive linear clause OR a subset of earlier marked outputs. Indeed, between private products only the old degree-one additions and the three cut insertions remain. Induction through those additions gives this form. An earlier unmarked product cannot be added to an old clause without another deleted addition. Carrier inputs have already been erased. Making each marked output available immediately after its creation enlarges the possible circuits and includes the original one.

Temporarily replace each high-input occurrence by an independent symbol of weight equal to the minimum degree of its low operand. The exact quadratic frame has weighted degree two. A weight-two symbol has only a scalar context; a weight-one symbol has a linear context or a second weight-one symbol. The square of one symbol cannot occur, since substituting its nonzero old linear low part would give a square in the squarefree quadratic polynomial, and positive coefficients cannot cancel it. Thus, on the restriction, the output has the form
\begin{equation}\label{eq:threevisibleform}
 \bigvee_{i=1}^3 g_iL_i\ \vee\!
 \bigvee_{1\leq i<j\leq3:\,\delta_{ij}=1}g_ig_j,
 \qquad \delta_{ij}\in\{0,1\},
\end{equation}
where $g_i$ are the marked products, repetitions allowed, and each $L_i$ is an affine positive clause. Allowing all pair terms is harmless even when particular weights would forbid some of them.

Every nonzero nonconstant $L_i$ contains an isolated variable or both endpoints of a matching edge. In fact, such a direct context comes from perturbing an assigned center $h\in Z$ and contains the nonempty set $S=N_H(h)\cap T$. If the context contained no isolate and no complete matching pair, then $S$ would meet exactly $|S|$ different matching edges. Replace $h$ in $Z$ by $S$, and add one endpoint of each remaining missing edge. This covers $H$: old edges from $h$ into $T$ meet $S$, and old edges between $h$ and another center meet the other center. It also covers the four missing edges, with at most $|Z|-1+|S|+4-|S|=|Z|+3$ vertices. This contradicts $d=4$.

If the three cuts select at most two distinct product nodes, collect the terms in \eqref{eq:threevisibleform} with the same high value. A repeated square is Booleanly that value and hence has constant context one. Unions of permitted contexts are still permitted. The resulting function belongs to the two-high-signal class excluded in Proposition~\ref{prop:fiveproductseven}, whose proof permits fewer than five products and does not require the restricted signals to be nonzero or distinct. Thus it remains to consider three distinct marked products.

Every original private product is ancestral to a marked product. Its first entry into the selected subgraph must be along the high-input edge of a visible cut: entry through a product, an equal-degree sum, or a cut's low input would make the private node selected already. The last of the four topologically ordered products must therefore itself be marked. Consequently the three marked index sets are exactly
\[
 \{1,2,4\},\qquad \{1,3,4\},\qquad \{2,3,4\}.
\]
Keep these original indices even if two restricted values coincide or vanish.

For completeness, the finite obstruction for these three classes is specified as follows. Identify all isolated variables with one variable $z$, leaving the target matching unchanged. In each product operand select exactly one of two kinds of input: an arbitrary old linear clause in the eight endpoints and $z$, optionally joined by any earlier marked outputs; or one earlier private-product output alone. The empty clause is allowed. Select each affine context by its nine variable-membership bits and a constant-one bit. Require a nonzero nonconstant context to contain $z$ or a full matching pair, and select the three pair terms independently. No essentiality, nonvanishing, or distinction of restricted truth functions is imposed.

The circuit must be true on each of the four edge assignments with $z=0$ and false on each of the sixteen maximal independent endpoint transversals with $z=1$. These twenty conditions imply equality on every assignment by monotonicity: every edge-containing assignment dominates an edge test, and every independent assignment is below a tested transversal. Exact Tseitin equivalences express the specified evaluations. The resulting CNFs have, respectively,
\[
 (3342,11778),\qquad(3300,11618),\qquad(3258,11458)
\]
variables and clauses. Complete DRUP proofs for them contain $196\,938$, $236\,242$, and $286\,925$ proof lines. Independent regeneration reproduced every CNF byte for byte. DRAT-trim verified each refutation in RUP-only mode with zero RAT steps, in $6.20$, $7.47$, and $10.62$ seconds. The encoder \emph{Visible-cut obstruction}, compressed refutations, hashes, regeneration checks and checker logs are available through the workbench. Reverse unit propagation establishes each added clause from earlier clauses and ends in the empty clause, as in Lemma~\ref{lem:threeprivatefinite}. These three refutations exclude every remaining class and prove the proposition.
\end{proof}

The remaining visibility cases have one or two visible cuts, considered next.

\subsection{One visible cut in the seven-deletion case}
The case of a single visible cut reduces to a small Boolean circuit with unrestricted sharing. The following exact finite bound treats that sharing explicitly.

\begin{lemma}[Four conjunctions and two disjunctions]\label{lem:fourandtwoor}
A Boolean circuit with at most four binary AND gates and at most two binary OR gates, taking arbitrary positive clauses as boundary inputs, has a positive CNF with at most twelve clauses. Reuse of gate values is unrestricted, and the boundary clauses may be zero or one. The bound twelve is attained. Conjoining one further affine positive clause gives at most thirteen clauses.
\end{lemma}
\begin{proof}[Proof by complete finite enumeration]
First remove gates not ancestral to the designated output. A circuit
with no gate is already one clause. Otherwise append idempotent gates, using the previous output twice, until there are exactly four AND and two OR gates. This preserves the Boolean function. Place the output last in a topological ordering, and number the six gates from zero through five. At gate $i$, an operand is a previous gate or one free clause. Replace each free-clause input occurrence by a distinct formal atom. Recovering the original circuit means substituting a positive clause for each atom, including any repeated or overlapping old inputs. Such substitution maps one formal CNF clause to one positive clause; false atoms may make a clause empty, and true atoms may remove a clause. Thus it cannot increase the number of clauses.

Represent a positive CNF by the inclusion-minimal sets of atoms occurring in its clauses. For clause families $\mathcal A,\mathcal B$, let $\min(\mathcal A)$ mean deletion of duplicate sets and of every strict superset of another set. The exact gate recurrences are
\[
 \mathcal A\mathbin{\mathrm{AND}}\mathcal B
   =\min(\mathcal A\cup\mathcal B),\qquad
 \mathcal A\mathbin{\mathrm{OR}}\mathcal B
   =\min\{A\cup B:A\in\mathcal A,\ B\in\mathcal B\}.
\]
These are distributivity and absorption. A shared gate reuses its existing family, so shared values are not replaced by independent copies.

There are $i+1$ input-source choices at gate $i$: its $i$ predecessors and the free-clause alternative. Since AND and OR are commutative, choose an unordered pair with repetition. If both sources are free clauses, use two different fresh atoms. Hence there are $\binom{i+2}{2}$ possibilities at gate $i$, and $\binom62=15$ operation patterns. The complete count is
\[
 15\prod_{i=0}^5\binom{i+2}{2}
 =15(1\cdot3\cdot6\cdot10\cdot15\cdot21)=850\,500.
\]
Every earlier-gate choice, repeated operand, free-clause occurrence and unused node is included; no semantic pruning is used.

Two separately written enumerators evaluate the displayed set recurrences on this entire finite class. The maxima, grouped by the positions of the two OR gates, are
\[
\begin{array}{c|rrrrrrrr}
\text{positions}&01&02&03&04&05&12&13&14\\\hline
\text{maximum}&5&5&6&7&9&6&7&8
\end{array}
\qquad
\begin{array}{c|rrrrrrr}
\text{positions}&15&23&24&25&34&35&45\\\hline
\text{maximum}&9&7&9&12&9&12&12.
\end{array}
\]
The second enumerator and its universal reduction were completed from the statement alone before comparison with the first. Their exact source and output hashes, all counts, and maximizing circuits are recorded with the program \emph{Shared-circuit clause enumeration} in the workbench. The largest number is twelve.

For an instructive sharp example, take seven distinct atoms and the function
\[
 (a_1\wedge a_2)\vee(b_1\wedge b_2)
       \vee(c_1\wedge c_2\wedge c_3).
\]
It has four AND gates and two OR gates. Its clauses are exactly
$a_i\vee b_j\vee c_k$, for $1\le i,j\le2$ and $1\le k\le3$; all twelve are indispensable. Finally, an additional conjunct appends at most one clause.
\end{proof}

\begin{proposition}[One visible cut]\label{prop:onevisible}
Under the hypotheses of Proposition~\ref{prop:threevisible}, exactly one of the three deleted additions cannot be a visible cut.
\end{proposition}
\begin{proof}
Apply full-carrier erasure to the original circuit. No private gate can disappear while preserving the exact quadratic extraction, since the resulting circuit would have at most six deletions in the excluded optimal-frame case. Every private node reaches the output and first enters the selected subgraph along the high-input edge of the unique visible cut. Therefore every private gate is ancestral to its high operand $g$. By acyclicity, $g$ does not depend on the visible cut's output.

The high-input circuit consequently has four products and two private additions, and its boundary values are old linear clauses. Full-carrier erasure excludes old quadratic inputs. Booleanly its restricted value has the form covered by Lemma~\ref{lem:fourandtwoor}. All original nodes remain allowed even if some become zero or redundant after restriction.

The independent-symbol quadratic-frame identity gives output $H\vee gL$, with one affine positive direct context $L$. The high symbol cannot enter both factors of a frame product: substituting its nonzero old linear low part would produce a square in the quadratic output. A degree-two insertion has only a scalar context. These observations cover both possible cut degrees.

Set the assigned minimum cover $Z$ to zero. The quadratic output $H$ disappears, and the target is a four-edge matching with possible isolated variables. Yet $gL$ has a positive CNF with at most thirteen clauses. Every clause true on each of the four individual matching edges must contain an endpoint of each edge. It can therefore reject at most one of the sixteen maximal independent endpoint transversals, with all isolated variables set to one. Constant-one clauses reject none, and an empty clause would reject the required edge assignments. At least sixteen clauses are necessary, a contradiction.
\end{proof}

Only the two-visible-cut branch remains in this profile.

\subsection{Two visible cuts and one private addition}
The remaining addition can either be absorbed into the earlier exposed signals or join the two private products that are not exposed. This distinction gives the last finite obstruction.

\begin{proposition}[Two visible cuts]\label{prop:twovisiblefour}
Under the hypotheses of Proposition~\ref{prop:threevisible}, exactly two of the three deleted additions cannot be visible cuts.
\end{proposition}
\begin{proof}
Apply full-carrier erasure before restriction, and let $q$ be the unique deleted addition that is not visible. The other additions are visible cuts. As before, the exclusion of at most six deletions implies that every original private gate remains and reaches one of their high operands. Except for $q$, a nonproduct operand in this complementary circuit is an old linear clause, possibly with earlier visible high signals inserted. Full quadratic carriers have been erased from its inputs.

First suppose $q$ itself is a cut's high operand. It has minimum degree at least two, so both its operands do also. Neither operand can be old linear or a degree-one visible cut; a degree-two visible cut is a protected carrier and is also excluded. Thus $q=u\vee v$ for two private products, possibly identical. Replace every Boolean occurrence of $q$ by $u\vee v$ and expose $u,v$ at their creation times. The second visible high operand is another private product or $q$ itself. There are now at most three exposed product outputs. Their direct contexts are the original permitted contexts, duplicated when needed, and pair terms expand to pair terms; repeated squares are absorbed into constant-one contexts. A private operand previously equal to $q$ becomes the permitted linear-mode operand with zero old part and exposed signals $u,v$. Thus the function belongs to the three-high or two-high relaxation already excluded. Every product is ancestral to a designated high, so the last product remains designated.

We may therefore assume that both visible highs are private products $g,f$. If both operands of $q$ are old or perturbed linear clauses, their union is another permitted linear-mode operand, giving the two-high relaxation. If its operands are a private product $u$ and a linear clause, expose $u$ as a third signal, with zero direct context and no new output pair terms. The value of $q$ is then a permitted linear-mode operand. The same argument applies when one of two product operands is already designated: include that signal in the linear part and expose the other one if necessary. Repeated identical operands need no separate exposure. These reductions enlarge the Boolean class; they do not assert that the original arithmetic circuit can lose an addition while preserving its extraction.

If $g=f$, expose both product operands of $q$ as two further signals, with zero direct contexts and no new output pair terms. There are at most three designated products, and $q$ is their permitted linear-mode sum. Every product still reaches the designated $g$, so the three-high relaxation applies. Thus only $q=u\vee v$ with two distinct undesignated private products and distinct $g,f$ remains. The four products are exactly $u,v,g,f$. Every product is ancestral to $g$ or $f$, so the last topological product is designated. Numbering the products from one to four, their designated sets are
\[
 \{1,4\},\qquad\{2,4\},\qquad\{3,4\}.
\]
The other two products are the operands of $q$. Allow $q$ immediately after both have been created. Each product operand is an old linear clause OR any subset of preceding designated highs, a previous product alone, or $q$ alone when available. The restricted output is
\[
 gL\vee fM\vee\delta gf,
 \qquad\delta\in\{0,1\},
\]
with the same allowed affine contexts as in Proposition~\ref{prop:threevisible}.

The three exceptional templates are again finite over nine inputs, after identifying all isolates with one variable. Select the stated input modes, all old variable memberships, all earlier designated-high memberships, both admissible affine contexts, and the pair term independently. Empty old clauses and zero contexts are allowed. No restricted gate is required to be nonzero, distinct or essential. Exact Tseitin equivalences impose truth on the four edge assignments and falsity on the sixteen maximal independent transversals; monotonicity makes these tests complete.

The three CNFs have $(3028,10563)$, $(2986,10403)$ and $(2986,10409)$ variables and clauses. Their complete DRUP proofs have $47\,892$, $54\,127$ and $93\,835$ lines. Independent regeneration reproduced the CNFs exactly, and independent RUP-only checking verified the refutations, with zero RAT steps, in $0.97$, $1.07$ and $2.08$ seconds. The complete data and hashes accompany the visible-cut obstruction. Hence all three exceptional classes are impossible, completing the proof.
\end{proof}

\begin{theorem}[The optimal seven-deletion case]\label{thm:sevendeletion}
Suppose exactly four quadratic witnesses are missing, $d=4$, and the extracted quadratic circuit is optimal. Then seven deletions are impossible. Consequently a circuit below the star bound with four missing witnesses must satisfy
\[
\begin{array}{c|l}
d&(D,e)\\\hline
2&(5,0)\\
3&(5,0),(5,1),(6,0)\\
4&(6,1).
\end{array}
\]
In the last row exactly four products and two additions are deleted.
\end{theorem}
\begin{proof}
Proposition~\ref{prop:fiveproductseven} leaves only four products and three additions. At least one addition is a visible cut; otherwise the selected degree-two output would equal the full output and no quadratic witness would be missing. The one-, two- and three-visible-cut cases are excluded by Propositions~\ref{prop:onevisible}, \ref{prop:twovisiblefour} and~\ref{prop:threevisible}. The displayed alternatives follow from the earlier four-witness reduction and the six-deletion exclusion.
\end{proof}

\subsection{Four matching witnesses over a quadratic circuit with excess one}
Full-graph safety provides the missing restriction on insertion contexts when the quadratic circuit is larger than its minimum. We first isolate the functional equality case in the four-product argument.

\begin{lemma}[The four-product matching alternative]\label{lem:fourproductalternative}
Suppose a positive circuit has nonempty squarefree quadratic support, six deleted gates consisting of four products and two additions, and eight vertices on which its Boolean output is a four-edge matching while its quadratic part vanishes. After full-carrier erasure, the only possible restricted computation has two visible cuts, zero independent-symbol pair coefficient, and the form
\[
 gL\vee fM,
\]
where $g$ and $f$ are independent products of three old linear clauses. Each safe quartic witnesses exactly two matching edges, and these two pairs are disjoint. Its four restricted clauses are exactly the four two-element transversals of its witnessed pair.
\end{lemma}
\begin{proof}
The following is the functional portion of \citet[proof of Lemma~5.35]{Companion}, before its use of quadratic optimality. We give the argument to specify the equality case needed here. Full-carrier erasure cannot remove a private gate. Otherwise the exact quadratic extraction would be preserved with at most five deletions. The four-product obstruction still forces at least four private products, and a missing witness forces at least one visible cut. The only possible profile would then have four products and one addition. Its unique high input is a multiplication DAG with at most five old linear boundary clauses, and the output context adds at most one clause. The resulting matching would have a CNF with at most six clauses, contrary to the required sixteen. This essentiality concerns the original circuit, before restriction.

If there is just one visible cut, all private gates feed its high operand. It has four products and at most one private addition over old linear clauses. Lemma~\ref{lem:fourandtwoor}, followed by the affine direct context, gives at most thirteen CNF clauses. The matching requires sixteen, so both additions are visible cuts. Their high inputs are private products $g,f$, and the exact restricted frame is
\[
 F=gL\vee fM\vee\delta gf,\qquad\delta\in\{0,1\}.
\]
The contexts are arbitrary affine positive clauses at this stage. Let $k$ count the private products ancestral to the first high $g$ in the original erased circuit. Every other private product feeds $f$. Keep the four nodes even if some restricted values vanish or coincide.

If $k=4$, the second high is an ancestor of $g$ or equals it. Write $g=fA$ Booleanly. Their total number of distinct old factors is at most five, and
\[
 F=f(AL\vee M\vee\delta A)
\]
has at most six clauses after distribution and absorption.

Suppose $k=3$. Then $g$ has at most $b=4$ old factors and only one product remains for $f$. If both inputs of $f$ are private ancestors of $g$, its factors are among those of $g$, so the preceding factorization applies. If exactly one input is a proper private ancestor $u$, write $g=uW$. The other input gives $f=uA$ or $uA\vee g$, including zero old parts. If $u$ has $r\ge1$ factors and $W$ has $t$ remaining factors, then $r+t\le4$. Factoring $u$ leaves, respectively,
\[
 WL\vee AM\vee\delta WA,
 \qquad W(L\vee M\vee\delta)\vee AM.
\]
The bracket has at most $2(t+1)$ clauses; with $u$ this gives at most $r+2(t+1)\le9$. For the first bracket with $\delta=1$, the identity
\[
 WL\vee AM\vee WA=(W\vee A)(W\vee M)(A\vee L)
\]
gives the slightly smaller count $2t+1$.

Without a proper private-ancestor operand, $f$ is $AB$, $AB\vee gB$, or $AB\vee g$. Direct distribution gives the following bounds:
\[
\begin{array}{c|cc}
f&\delta=0&\delta=1\\\hline
AB&3(b+1)\le15&3b+2\le14\\
AB\vee gB&3b+2\le14&3b+1\le13\\
AB\vee g&3b+1\le13&3b\le12.
\end{array}
\]
For example, in the middle row with $\delta=0$ the identity is
$F=(g\vee ABM)(L\vee BM)$; with $\delta=1$ it is
$F=(g\vee ABM)(L\vee B)$. In the last row absorption gives
$F=(g\vee ABM)(L\vee M)$ or $F=g\vee ABM$.
All these counts are below sixteen.

Suppose $k=2$. Write $u=UV$ and $g=uW$, allowing $W=1$ when the second product repeats its input. If the two-product network for $f$ reads $u$ directly, that factor persists to $f$. Under it a $g$-affected boundary form $A\vee g$ becomes $A\vee W$. Thus $f=uK$, with $K$ a product of at most two linear clauses, and
\[
 F=u(WL\vee KM\vee\delta WK).
\]
The bracket has at most six clauses for $\delta=0$, and at most five for $\delta=1$ by
$(W\vee K)(W\vee M)(K\vee L)$. Including $U,V$ gives at most eight.

Otherwise $g$ is a product of at most $b=3$ old clauses, and $f$ is a product of at most $t=3$ old or $g$-affected clauses. This leaf count includes sharing: a connected two-product multiplication graph has at most three distinct boundary inputs. If none is affected, write $f=K$. For $\delta=0$ the clause count is at most $(b+1)(t+1)\le16$; for $\delta=1$ the identity
\[
 gL\vee KM\vee gK=(g\vee K)(g\vee M)(K\vee L)
\]
gives $bt+b+t\le15$ clauses.

If at least one factor is affected, distributivity and absorption write
$f=H_0K_0\vee gH_0$, where $H_0$ has at most $h\le2$ old factors and $H_0K_0$ has at most $t\le3$. Then
\[
 F=(g\vee H_0K_0M)(L\vee H_0M)\quad(\delta=0),
\]
and the second factor becomes $L\vee H_0$ for $\delta=1$. The counts are at most $b(t+1)+h+1\le15$ and $b(t+1)+h\le14$.

Finally, if $k=1$, $g$ has at most $b=2$ factors and the remaining three products give at most $t=4$ boundary factors for $f$. If none is affected, the preceding counts are at most fifteen and fourteen. Otherwise $h\le3$ and the counts are at most $b(t+1)+h+1\le14$ and $b(t+1)+h\le13$. A directly used $g$ is included by setting the old part of its affected clause to zero. Repeated factors and restricted zero values only decrease these bounds.

The only count that can reach sixteen is therefore $k=2$, $\delta=0$, with $g,f$ independent products of exactly three old clauses. Each summand $gL$ or $fM$ is safe and has at most four clauses. A safe CNF true on $s$ matching edges needs at least $2^s$ clauses, by the independent-transversal argument; hence each summand witnesses at most two. They must each witness two, with disjoint pairs, to cover all four edges.

For a safe quartic witnessing two edges, each clause meets both edge assignments. The four independent transversals must all be rejected, and one clause can reject at most one. Equality therefore forces the four distinct complementary two-element clauses. No clause can contain an outside variable: adjoining that variable to its uniquely rejected transversal would make all four clauses true on an independent assignment. A clause cannot be constant either. These are exactly the four asserted transversals.
\end{proof}

\begin{proposition}[Full insertion contexts force a small cover]\label{prop:sixmatchingcover}
Let $G$ have no four-cycle, and let a positive variable-input binary circuit compute $P\in\Sup(G)$ with nonempty quadratic support $H$. Suppose exactly four edges lack quadratic witnesses and their eight endpoints induce precisely those four edges in $G$. If exactly six gates are deleted, four of them products, and $K$ products remain in the quadratic circuit, then
\[
 \tau(G)\le K+2.
\]
Consequently this circuit has at least $|E(G)|+\tau(G)-1$ gates.
\end{proposition}
\begin{proof}
Apply full-carrier erasure and restrict to the eight endpoints. By Lemma~\ref{lem:fourproductalternative}, the sole possible computation has two visible cuts and the independent-symbol form $P_2+gD_g+fD_f$, with zero pair coefficient. Here $D_g,D_f$ are the \emph{full} old direct contexts, before restriction. Relabel the missing edges
\[
 aa',\quad bb',\quad cc',\quad dd'
\]
so that the restricted contexts are $a\vee b$ and $c\vee d$. Quartic equality gives
\[
 g(\{a',b'\})=1,\qquad f(\{c',d'\})=1.
\]
All these statements concern Boolean supports; positive coefficients do not change them.

Let $A_g$ be the nonzero old linear low part of the $g$ cut. The product $A_gD_g$ occurs in the squarefree quadratic frame. Since $D_g$ contains two distinct variables $a,b$, four-cycle-freeness forces $A_g$ to have singleton support, say $\{h\}$: two distinct variables in each disjoint factor would give a four-cycle. Squarefreeness makes the factor supports disjoint. The edges $ha,hb$ belong to $H$, so $h$ lies outside the eight endpoints, which contain no $H$-edge.

Suppose $D_g$ contains an additional vertex $z$. Its restriction has only $a,b$, so $z$ also lies outside the eight endpoints. The edge $hz$ lies in $H$, and $z\ne h$ by squarefreeness. On the assignment $\{a',b',z\}$ the positive contribution $gD_g$ is true. Since $a'b'$ is not an edge, correctness forces $za'$ or $zb'$ to be an edge of $G$. In the first case $h,a,a',z$ form a four-cycle, and in the second $h,b,b',z$ do. All four vertices are distinct. This contradiction proves
\[
 \operatorname{supp}(D_g)=\{a,b\}.
\]
The identical argument gives $\operatorname{supp}(D_f)=\{c,d\}$. The full contexts, not merely their restrictions, are disjoint two-element sets.

Call a quadratic product affected by a cut if its independent degree-one symbol reaches one of the product's two linear factors. Each cut affects at least one product, because its direct context is nonzero and linear. It cannot affect both factors of one product: substituting its nonzero old low part would produce a square. If a product is affected by the $g$ cut, its old opposite factor contributes positively to $D_g$, so its support is contained in $\{a,b\}$. Thus these two vertices cover every edge of that product star, including edges from other old summands in the perturbed factor.

No product is affected by both cuts. If both symbols enter the same factor, the nonempty opposite factor contributes to both disjoint contexts. If they enter opposite factors, they produce a positive independent-symbol pair coefficient, which has been excluded. Both conclusions use positivity, so there is no cancellation. They allow shared inputs, shared intermediate values and nested cut locations.

The two disjoint affected sets therefore contain at least two quadratic products. Their edges are covered by $S=\{a,b,c,d\}$, which also covers all four missing edges. For each unaffected product choose one star center. There are at most $K-2$ such products, and their centers together with $S$ cover $G$. Hence $\tau(G)\le K-2+4=K+2$.

We justify the intermediate count $|C_2|\ge |E(H)|+K-1$ without an optimality assumption. Every quadratic product is a star: both factors are linear, their supports are disjoint by squarefreeness, and two variables in each factor would create a four-cycle. Choose its singleton factor as its center, choosing either endpoint for a one-edge product. Suppress linear additions with singleton support; this preserves all $K$ products and the positive edge support, and only decreases size.

A remaining linear wire has at least two variables in its support. It cannot feed products with distinct centers, since those two centers and any two of its variables would form a four-cycle. Thus these linear gates partition by their unique product center. For a center $v$, let $k_v$ be the number of its products, $q_v$ the number of distinct variables in the union of their leaf factors, and $a_v$ the number of linear additions assigned to it. The binary input graph for these factors has at most $k_v$ connected components, each containing one of the factor outputs. Counting its vertices and edges gives $a_v\ge q_v-k_v$. This also includes single-variable factor outputs and repeated use of one factor wire. Every edge of $H$ occurs in at least one chosen star, so $\sum_vq_v\ge|E(H)|$ and $\sum_vk_v=K$. The linear part therefore has at least $|E(H)|-K$ additions. Above the $K$ distinct product nodes the output network consists only of additions and has at least $K-1$ gates by the same connected-graph count. Adding the $K$ products proves the claimed inequality.

Since $|E(G)|=|E(H)|+4$ and six gates were deleted,
\[
 |C|\ge |E(H)|+K+5
       =|E(G)|+K+1
       \ge |E(G)|+\tau(G)-1.
\]
\end{proof}

\begin{theorem}[Maximal cover gain with four missing witnesses]\label{thm:fourmaxgain}
Let $G$ be a finite simple graph without a four-cycle. If a positive variable-input binary circuit computes a polynomial in $\Sup(G)$ with exactly four missing quadratic edge witnesses and nonempty quadratic support $H$, and if $\tau(G)-\tau(H)=4$, then its size is at least $|E(G)|+\tau(G)-1$.
Thus the remaining possible pairs for a four-witness counterexample are
\[
\begin{array}{c|l}
d&(D,e)\\\hline
2&(5,0)\\
3&(5,0),(5,1),(6,0).
\end{array}
\]
\end{theorem}
\begin{proof}
The previous reductions leave only $(D,e)=(6,1)$ when $d=4$. Maximal cover gain forces the missing edges to be an induced four-edge matching with no old edge on their endpoints. Exactly four products and two additions are deleted. Proposition~\ref{prop:sixmatchingcover} excludes a saving in this final case. The displayed smaller-gain alternatives are those of the preceding reduction.
\end{proof}

\subsection{Five deletions for two missing-edge configurations}
The cover-gain-three cases depend on the shape of the missing edges. Two shapes permit an additional exclusion.

\begin{proposition}[A path or triangle and a disjoint edge]\label{prop:pathfive}
Let $G$ have no four-cycle, and let a positive variable-input binary circuit compute $P\in\Sup(G)$ with nonempty quadratic support $H$. Suppose the graph formed by its four missing edges is $P_4\sqcup K_2$ or $K_3\sqcup K_2$, where $P_4$ is a path with three edges. If $d=3$ and the extracted quadratic circuit is optimal, then five deletions are impossible.
\end{proposition}
\begin{proof}
Let $J$ be the missing-edge graph, $t=\tau(H)$, and $Z$ the assigned minimum cover of the optimal quadratic frame. Both possible $J$ have cover number three, and deleting any one of their vertices lowers that number to two. If such a vertex belonged to $Z$, adjoining a cover of the remaining missing graph would cover $G$ with $t+2$ vertices. Hence $Z$ avoids every missing endpoint. On $T=V(G)\setminus Z$, the graph is precisely $J$ with possible isolated vertices.

Apply full-carrier erasure before restriction. No private gate disappears, by the earlier five-deletion lower bound for four missing witnesses. Every nonzero nonconstant restricted direct context of a visible cut is not contained in a minimum cover of $J$. Indeed, it perturbs a literal assigned center $h\in Z$ and contains $N_H(h)\cap T$. If it lay in a minimum $J$-cover $U$, replacing $h$ in $Z$ by $U$ would cover $H$ and all missing edges with $t+2$ vertices. No missing edge is incident to $h$, since $Z$ avoids their endpoints. Zero and scalar contexts remain allowed. A union of permitted nonzero linear contexts is still not contained in any minimum $J$-cover.

A second restriction comes directly from four-cycle-freeness. If $u,w,v$ form a two-edge path of missing edges, no old linear boundary clause may contain both $u$ and $v$. Such a surviving clause has a quadratic product above it and therefore a common old neighbor $h$ of $u,v$. Squarefreeness excludes $h=u,v$; $h=w$ would make a missing edge old; and any other $h$ produces the four-cycle $u,w,v,h$. The same restriction holds for one individual nonconstant visible-cut context $D$, because its nonzero old low part $A$ satisfies $AD\subseteq P_2$ in support. It need not hold for a union of contexts from different cuts.

Identify all isolated vertices with one variable $z$. This preserves the target, the minimum-cover context condition, and the missing-path restrictions on the genuine endpoints. We now consider the number of deleted products.

With at most two products, full-carrier expansion gives the general two-product form
\[
 g_i=A_i(x,g_0,\ldots,g_{i-1})B_i(x,g_0,\ldots,g_{i-1}),\qquad
 F=\bigvee_i g_iL_i\vee\!\bigvee_{i<j:\delta_{ij}=1}g_ig_j.
\]
The operand clauses are unrestricted positive disjunctions, and the affine contexts obey only the minimum-cover condition. A direct context is a union of visible-cut contexts; diagonal terms become constant-one contexts by Boolean idempotence. In this class we impose no missing-path restriction on the operands or the merged contexts. Private additions can create new linear sums, and different cuts can have different old neighbors. The finite certificate below excludes this larger class for both targets.

Suppose there are three products and two additions. If both additions are visible cuts selecting the same high product, all three products are ancestral to it, and it is a conjunction of at most four old clauses. One merged affine context gives at most five CNF clauses, whereas either target requires six. If the two highs are distinct, no private addition remains. Each operand is an old linear clause, possibly with inserted highs, or an earlier product alone. We enlarge this class by allowing every earlier product in any operand disjunction. Its old variable part satisfies the missing-path restriction, and each of the two individual direct contexts satisfies both restrictions. Only two product indices can appear directly or in output pair terms. The second finite class below excludes this possibility.

If only one addition is visible, the high-input computation is a four-node Boolean DAG with three AND gates and one OR gate, over old linear boundary clauses. Each old clause obeys the missing-path restriction. The private OR itself may create a forbidden pair, so it is represented as an explicit gate with unrestricted output. Each input is an old clause or an earlier node, and the final affine context obeys both restrictions. The third finite class below excludes this case, including every operation order and all sharing.

Finally suppose four products are deleted. The sole deleted addition would contradict Theorem~\ref{thm:purecut}, since $d=3$ requires $D\ge2^3-1=7$. At least one deleted addition is necessary, so these product counts exhaust the possibilities.

Here are the complete finite verifications used in the first three cases. For the first two classes, choose each variable and each permitted earlier product in an operand by a Boolean selector, and select affine contexts and pair terms similarly. The minimum-cover condition says that a nonzero nonconstant context has a selected variable outside each minimum cover. In the second class, prohibit each missing-path pair in old variable parts and individual contexts; three activity bits dominate every context and pair selector, with at most two bits true. In the third class, one-hot operand modes choose an old clause or an earlier node, and operation bits select exactly one OR among four gates. Its old clauses and final context obey the stated restrictions.

Exact Tseitin equivalences impose truth on the four edge assignments with $z=0$ and falsity on the six maximal independent assignments with $z=1$. Monotonicity makes these ten tests complete. Empty clauses, zero or constant-one contexts, repeated operands, unused nodes and restricted coincidences are permitted. The resulting counts are
\[
\begin{array}{c|l|r|r|r}
J&\text{class}&\text{variables}&\text{clauses}&\text{proof lines}\\\hline
P_4\sqcup K_2&\text{two products}&632&2157&493\\
P_4\sqcup K_2&\text{two visible cuts}&992&3442&2210\\
P_4\sqcup K_2&\text{one private OR}&1130&3838&1678\\
K_3\sqcup K_2&\text{two products}&566&1913&524\\
K_3\sqcup K_2&\text{two visible cuts}&893&3085&2873\\
K_3\sqcup K_2&\text{one private OR}&1031&3486&1093.
\end{array}
\]
All six CNFs were independently regenerated byte for byte. Complete DRUP refutations were independently verified in RUP-only mode with zero RAT steps; the checks took between $0.06$ and $0.27$ seconds each. The program \emph{Five-deletion obstruction}, exact source bindings, compressed proofs and checker logs accompany the workbench. The certified classes contain all cases above, completing the proof.
\end{proof}

At cover gain three with five deletions and an optimal frame, the missing graph therefore cannot be $P_4\sqcup K_2$ or $K_3\sqcup K_2$. The shape $P_3\sqcup K_2\sqcup K_2$, and configurations in which a minimum old cover meets a missing endpoint, require further arguments. A center exchange in the latter case can uncover a missing edge incident to the removed center.

\subsection{Six deletions for a path and a disjoint edge}
The path configuration permits a further exclusion when triangles are also
absent. The additional restriction concerns how distinct inserted values
can occur on opposite sides of a quadratic product.

\begin{proposition}[An optimal frame for a missing path and edge]\label{prop:pathsix}
Let $G$ have girth at least five, and let a positive variable-input binary
circuit compute $P\in\Sup(G)$ with nonempty quadratic support $H$.
Suppose its four missing edges form $P_4\sqcup K_2$, the cover gain is
three, and the extracted quadratic circuit is optimal. Then six deletions
are impossible. Consequently every circuit with this missing-edge shape,
cover gain three and an optimal quadratic frame satisfies the full star
lower bound.
\end{proposition}
\begin{proof}
Write the missing edges as $ab,bc,cd,ef$ and call their graph $J$.
The assigned minimum $H$-cover $Z$ avoids all six endpoints, as in
Proposition~\ref{prop:pathfive}. On $T=V(G)\setminus Z$, the graph is
therefore $J$ with possible isolated vertices. Full-carrier erasure loses
no private gate: a loss would leave at most five deletions with the same
optimal frame, already excluded by Proposition~\ref{prop:pathfive} and
the preceding four-witness deletion bound. This assertion concerns the
original circuit; restricted private values may be zero, redundant or
equal.

We use three restrictions from the quadratic frame. Every nonzero
nonconstant individual direct context on $T$ is not contained in a
minimum cover of $J$. Indeed, such a context contains all outside
neighbors of the literal assigned center it perturbs; replacing that
center by a minimum $J$-cover containing the context would cover $G$
with $\tau(H)+2$ vertices. A union of permitted contexts obeys the same
condition. Second, an old linear boundary clause and an individual
linear context avoid the pairs $\{a,c\}$ and $\{b,d\}$. A common old
neighbor for either pair would complete a four-cycle with its missing
two-edge path. The context assertion uses the positive old product
$AL\subseteq P_2$, where $A$ is the cut's nonzero old low part.
Third, triangle-freeness also prohibits adjacent vertices of $J$ in an
old clause or individual context: their common old neighbor would form
a triangle. We use this third restriction only in the case of three
distinct visible highs. Neither pair restriction is asserted for a
union of contexts from different cuts or for the output of a private
addition.

Identify all isolated vertices of $T$ with one variable $z$. The target
remains $J$, and the restrictions on its six genuine endpoints persist.
Zero and constant-one affine contexts are allowed throughout. With at
most two private products, the unrestricted two-product class in
Proposition~\ref{prop:pathfive} already gives a contradiction, regardless
of the number of private additions. With one deleted addition,
Theorem~\ref{thm:purecut} gives $D\ge 2^3-1=7$. At least one addition is
deleted. The remaining product/addition counts are $(4,2)$ and $(3,3)$.

\emph{Four products and two additions.}
If only one addition is visible, its high value is computed by a
five-node Boolean DAG with four AND gates and one OR gate over old
linear clauses. Every private node is ancestral to that high, because
it has no other exit to the selected quadratic subgraph. Its one
affine context obeys both the minimum-cover and missing-path
restrictions. This is the sixth finite class below.

If both additions are visible and select the same high $g$, no private
OR remains. The four-product multiplication DAG has at most five
distinct old boundary factors. Its merged context gives at most six
CNF clauses. The edge function of $J$ has six minimal vertex covers,
all of size three. Every positive CNF representing this function must
contain each of those covers as a clause: at its complementary maximal
independent assignment, some false clause is a cover contained in the
minimal cover, and hence equals it. A representation with at most six
clauses must therefore have exactly these clauses. A nonconstant
merged context cannot be one of them by the minimum-cover restriction;
a constant-one context leaves at most five clauses. Both alternatives
are impossible.

If the two highs are distinct, they are private products. Expose each
as soon as it is computed, which only enlarges its actual availability.
A private operand is an old linear clause joined with earlier exposed
highs, or one earlier private product alone. An unexposed product cannot
mix with another argument without a private addition, and none is
available. Each individual context obeys the two restrictions. Only
the two exposed indices may occur in direct output terms or pair
terms. This is the seventh finite class.

\emph{Three products and one visible cut.}
The high computation is a five-node DAG with three AND gates and two
private OR gates. Each input is an old clause or an earlier node, with
arbitrary sharing. Old clauses satisfy the missing-path restriction;
private OR outputs are unrestricted. The unique affine context
satisfies both restrictions. This is the first finite class.

\emph{Three products and two visible cuts.}
There is one private OR node $q$. If the highs coincide, all four
private nodes precede their common high, whose computation has three
ANDs and one OR. Its merged context obeys the minimum-cover condition
but need not avoid a missing-path pair. Diagonal terms are absorbed
into a constant-one context. This is the third finite class.

If the highs are distinct, each is a private product or $q$. An old
linear value or a degree-one visible cut is too small to be a high
operand. A degree-two visible cut is a protected carrier and is
unavailable by erasure. Thus the three products and $q$ form four
private nodes, at most two of which are designated highs. Each input
is an old clause joined with earlier designated highs, or an earlier
private node alone. The OR node is explicit, including when it joins
two old clauses; its output is not subjected to an old-clause pair
restriction. Each individual output context obeys both restrictions.
This is the second finite class.

\emph{Three products and three visible cuts, with repeated highs.}
No private addition remains, and every visible high is a private
product. If all highs coincide, at most four old factors and one
merged context give five clauses, fewer than the required six. If
there are exactly two distinct highs, one is used by two cuts and one
by a single cut. The former's merged context obeys only the
minimum-cover condition; the latter's individual context also avoids
the missing-path pairs. There is at most one mixed pair term, and
diagonal terms are absorbed into constant contexts. Unselected
products remain available alone but cannot mix with old arguments.
The fifth finite class permits at most two exposed output indices and
allows at most one context to violate the missing-path restriction.

\emph{Three products and three distinct visible highs.}
Here each private product is the high operand of one visible cut.
Let $S_i$ consist of the genuine missing endpoints $u$ on whose
singleton assignment the restricted high $g_i$ is true. The identified
isolate $z$ is excluded from $S_i$.

Call a cut center-affecting if its independent degree-one symbol
reaches the literal center factor of a quadratic star. Its old low
part is that literal center $h\in Z$, and its restricted direct
context contains $N_H(h)\cap T$, a nonempty set because $Z$ is minimum.
If $g_i$ were true on a singleton $u$, safety of its direct contribution
would put that context inside $N_J(u)$. Every neighborhood in
$P_4\sqcup K_2$ is contained in a minimum cover of $J$: the two internal
path neighborhoods are $\{a,c\}$ and $\{b,d\}$, and the other
neighborhoods are singletons extendable to such a cover. This
contradicts the context condition. Therefore
\[
 \text{a center-affecting cut has }S_i=\varnothing.             \tag{*}
\]

If the independent-symbol pair coefficient $\delta_{ij}$ is positive,
the two symbols enter opposite factors of a retained quadratic
product. One factor is its literal center, so at least one cut is
center-affecting. By $(\ast)$, $S_i$ and $S_j$ cannot both be nonempty.
A degree-two cut has scalar context and cannot create a quadratic
pair term.

There is also a common-neighbor restriction. Fix $j$ and put
$U_j=\bigcup_{i:\delta_{ij}>0}S_i$. If cut $j$ is never
center-affecting, every adjacent cut $i$ is center-affecting at the
corresponding product, and $U_j$ is empty by $(\ast)$. Otherwise the
old low part of cut $j$ is one fixed literal $h_j\in Z$ at every use.
Each pair term $g_i g_j$ then has a positive contribution $g_i h_j$
obtained by taking the old low value on the $j$ side. Safety on
$\{u,h_j\}$, for $u\in S_i$, forces the edge $uh_j$. It is an old edge
because every missing endpoint lies outside $Z$. Hence $U_j$ is
contained in one old neighborhood. It contains neither adjacent
vertices of $J$ nor the outer pair of a missing two-edge path, by
triangle- and four-cycle-freeness.

These two singleton restrictions require the distinct-cut hypothesis
just used. The fourth finite class imposes them only in the present
case. It otherwise allows all three products in later input
disjunctions, all individual affine contexts and all pair terms.
Old clauses and contexts avoid both the missing-path pairs and the
adjacent pairs, and contexts obey the minimum-cover condition.

\emph{Finite verification.}
In each class, Boolean selectors choose old arguments, permitted
earlier high values, contexts and pair terms. Exact bidirectional
Tseitin equivalences evaluate every selected argument and gate.
Truth on the four edge assignments with $z=0$ and falsity on the six
maximal independent assignments with $z=1$ are necessary and
sufficient for equality with the monotone target. The fourth class
also evaluates the six genuine singleton assignments to express its
two conditional restrictions.

For a one-cut DAG, one-hot modes select an old clause or a prior
node, and operation bits impose the exact AND and OR counts.
The class with one private OR and two distinct highs additionally
permits selected earlier highs in old linear inputs, and permits its
OR node to be selected as a high. Pure-product classes use activity
bits: an unselected prior product may be read alone, but mixing it
with another argument requires its activity bit. Context and pair
selectors imply the relevant activity bits. A relaxation bit permits
the one exceptional merged context in the repeated-high class.
All classes allow empty clauses, zero or constant-one contexts,
repeated source choices, unused nodes and restricted coincidences.

The seven exact refutations have the following sizes, in the order
of the classes specified above:
\[
\begin{array}{c|l|r|r|r}
 &\text{class}&\text{variables}&\text{clauses}&\text{proof lines}\\\hline
1&3\text{ ANDs},\ 2\text{ private ORs},\ 1\text{ cut}&1445&4932&19292\\
2&3\text{ ANDs},\ 1\text{ private OR},\ 2\text{ distinct highs}&1669&5856&19529\\
3&3\text{ ANDs},\ 1\text{ private OR},\ 1\text{ repeated high}&1130&3836&1678\\
4&3\text{ ANDs},\ 3\text{ distinct visible highs}&1595&5658&2912\\
5&3\text{ ANDs},\ 2\text{ distinct highs in }3\text{ cuts}&995&3491&1277\\
6&4\text{ ANDs},\ 1\text{ private OR},\ 1\text{ cut}&1445&4928&5200\\
7&4\text{ ANDs},\ 2\text{ distinct visible highs}&1389&4953&2932.
\end{array}
\]
Each CNF was independently regenerated byte for byte, and every
refutation was checked in RUP-only mode with zero RAT steps. The
checks took between $0.11$ and $0.42$ seconds. The program
\emph{Six-deletion path obstruction}, exact source bindings,
compressed refutations and checker logs accompany the workbench.
Together with the two-product and one-addition exclusions, these
classes exhaust the six-deletion cases. The last assertion follows
because the only possible counterexamples with an optimal frame and
cover gain three had five or six deletions.
\end{proof}

\section{The lowest component and the missing-edge charge}
The component operation used above costs no extra gates in the edge-support class, but does not by itself charge the missing edges. The underlying lower-envelope theorem is classical \citep[Theorem~2.4, p.~879]{JerrumSnir1982}; we spell out its application to the precise bidegree question.

\begin{proposition}\label{prop:componentzero}
Let $G=(X\sqcup Y,E)$ be bipartite, and let $P\in\mathbb N[X,Y]$ have the property that every monomial support contains an edge of $G$. A positive arithmetic circuit for $P$ gives a circuit for its bidegree-$(1,1)$ component with no increase in either additions or multiplications, provided a free zero output is permitted when that component is zero. If the component is nonzero, no constant input need be introduced.
\end{proposition}
\begin{proof}
For a nonzero polynomial $A$, let $\delta(A)$ be its least total degree and $\ell(A)$ its component in that degree. Positivity gives
\[
 \delta(AB)=\delta(A)+\delta(B),\qquad
 \ell(AB)=\ell(A)\ell(B).
\]
For addition, retain only the operand of smaller minimum degree when the two degrees differ; when they agree, add their lower components. Apply these rules in topological order, preserving every shared wire, and delete unused gates. No gate introduces more than one operation of its original type. If $P$ has a quadratic monomial, its minimum degree is two, and every quadratic monomial has one variable from each side. Therefore $\ell(P)=P_{1,1}$. If there is no quadratic monomial, the desired component is zero.
\end{proof}

Write $E'$ for the quadratic support. The extracted circuit computes a positive weighting of $H=(V,E')$, so its lower bound is $|E'|+\tau(H)-1$. To recover the conjectured bound for $G$, a separate argument must charge the deleted gates for both $|E\setminus E'|$ and $\tau(G)-\tau(H)$. A charge involving only additions is already impossible for the Petersen example in \citet[Section~6, example `A repeated-variable tradeoff on the Petersen graph']{OneSidedCompanion}.

Here is the exact computation. Put
\[
 u=x_0+x_6,\quad v=x_5+x_6,\quad w=x_0+x_9,
 \quad t=x_9+x_5,\quad p=x_7+x_3.
\]
The polynomial
\[
 P=x_4(w+x_3)+(x_1u+x_2)(p+x_1)
       +x_8(x_3+v)+(x_7+uvw)t
\]
has precisely the Petersen edges as its minimal monomial supports. Its
component of total degree two is
\[
 Q=x_4(w+x_3)+x_2(p+x_1)+x_8(x_3+v)+x_7t.
\]
This example concerns total degree, since the Petersen graph is not
bipartite. The missing quadratic edges are $01,05,16,69$; the other eleven are
\[
 04,34,49,12,23,27,38,58,68,57,79.
\]
Indeed the terms of $Q$ give those eleven edges. The remaining terms are $x_1u(p+x_1)$ and $uvwt$. Every monomial of the former contains $01$ or $16$, and its $x_1^2u$ terms witness both pairs. For the latter, Boolean absorption gives
\[
 (x_0\vee x_6)(x_5\vee x_6)(x_0\vee x_9)(x_9\vee x_5)
 =(x_0x_5)\vee(x_6x_9).
\]
To check this identity, the first two factors reduce to $x_6\vee x_0x_5$, the last two to $x_9\vee x_0x_5$, and their product is the displayed disjunction. Positivity then shows that every monomial support contains $05$ or $69$, and both occur. Thus the fifteen minimal supports are exactly the asserted graph; its usual outer-five-cycle and inner-star presentation, after relabelling, identifies it as the Petersen graph.


The displayed DAG for $P$ uses thirteen additions and seven products.
Its lower component $Q$ uses ten additions and four products after
unused wires are removed. Four edges have therefore lost quadratic
witnesses while only three additions are deleted. The total loss is six
gates, consistent with a possible joint charge. This example supplies no
counterexample to the total-size conjecture.

\section{Reductions at syntactic bidegree two}\label{sec:degree-two}

This section gives replacements that preserve shared wires and do not
increase the number of gates in two specified classes of circuits.
Let $G=(X\sqcup Y,E)$ be bipartite.  A positive polynomial is
\emph{safe} if every one of its monomials has support containing an
edge of $G$. Syntactic bidegree bounds the formal degree separately in
$X$ and $Y$: variable inputs have degrees $(1,0)$ or $(0,1)$, a product
adds the two input pairs, and an addition takes their componentwise
maximum. A useful gate with an addition-only route from that gate
to the output of a circuit for $f_G$ is safe.  In a variable-input
circuit of syntactic bidegree at most $(2,2)$, every useful gate of
bidegree $(2,2)$ has only additions above it.

\begin{lemma}[Four existing pure-side clauses]
Suppose a useful safe gate of a monotone circuit computing $f_G$
has Boolean value
\[
 g=A(X)B(X)C(Y)D(Y),
\]
where all four factors are existing nonempty linear disjunction wires.
If $G$ is $C_4$-free, then $g$ can be replaced by zero or by a single
star product using an existing factor wire and one variable, without
increasing the circuit size or changing its output function $f_G$.
Sharing of all four factor wires with the rest of the circuit is
permitted.
\end{lemma}
\begin{proof}
Write $A,B,C,D$ also for their variable supports, and put
$I=A\cap B$, $J=C\cap D$.  The gate accepts the edge assignment
$\{x,y\}$ exactly when $(x,y)\in I\times J$.  If either intersection
is empty, $g$ is false on every edge assignment. Replacing it by zero
leaves all these assignments unchanged and cannot create a false
positive, so the output remains $f_G$.  Otherwise safety gives $I\times J\subseteq E$.
The absence of a four-cycle implies $|I|=1$ or $|J|=1$.

Suppose $I=\{x\}$.  If both $C$ and $D$ contain a nonneighbor of $x$,
choose $y\in C\setminus N(x)$ and $z\in D\setminus N(x)$.
The gate accepts $\{x,y,z\}$, an independent set, which is impossible.
Hence $C\subseteq N(x)$ or $D\subseteq N(x)$.
Replace $g$ by $xC$ or $xD$, respectively, reusing the gate and the
specified factor wire.  This is a safe star and accepts every old
edge witness because $J\subseteq C,D$.  The case $|J|=1$ is symmetric.
The replacement is safe and still accepts every edge assignment
accepted by $g$. Every new output monomial passing through the replacement
therefore contains an edge, while all other output monomials are
unchanged. On each edge assignment the replacement cannot change a true
gate value to false, so monotonicity preserves every output edge witness.
Thus the new output is safe and still accepts every edge, hence equals
$f_G$. The pure-side linear factor wires cannot depend on $g$, whose
value requires variables of both sides, so this rewiring creates no
cycle. Zero propagation and pruning add no gates.
\end{proof}

The replacement uses the existing factor wires.  Merely expanding a sum
of bilinear products into many four-clause products would not preserve
the gate count.  The following simultaneous replacement avoids such
an expansion for a larger, explicitly restricted architecture.

\begin{theorem}[A shared bilinear two-level class]
Suppose a binary monotone circuit has the following form:
\begin{enumerate}
\item Its first product gates are $R_i=A_i(X)B_i(Y)$, with pure-side
linear disjunction operands, shared arbitrarily.
\item Its second product gates multiply two disjunctions of the $R_i$.
\item Its output is a disjunction of first and second product gates.
\end{enumerate}
All disjunction networks may have arbitrary fanout and sharing.
Suppose the circuit computes $f_G$ for a nonempty $C_4$-free bipartite
graph.  For every unsafe rectangle put
\[
 U_i=\{x\in A_i:B_i\subseteq N(x)\}.
\]
If $|U_i|\leq1$ for every unsafe rectangle, the circuit can be
replaced without growth by a quadratic circuit with edge support $E$.
In particular, $|B_i|\geq2$ for every unsafe rectangle is sufficient.  Consequently its size is at
least $|E|+\tau(G)-1$, using the established exact quadratic theorem (Theorem~\ref{thm:granted}).
The symmetric assertion holds with $X$ and $Y$ exchanged.
\end{theorem}
\begin{proof}
Keep every safe rectangle.  For an unsafe rectangle $R_i$, the
hypothesis gives at most one $x\in A_i$ with $B_i\subseteq N(x)$.
If that $x$ exists, replace $R_i$ by $xB_i$, using its existing
product gate and its existing $B_i$ wire.  If $U_i$ is empty, replace
it by zero.  Replace every second product gate by a disjunction gate.
All remaining first products are safe, so the resulting output is
safe; the replacements introduce no gate and do not change any shared
linear wire.

We show that every old edge witness remains.  A first product with
an addition-only route to the output was safe, and was kept.  If an
edge $xy$ was witnessed at a second product, choose first rectangles
$R_i=A_iB_i$ and $R_j=A_jB_j$ in its two operands which both accept
$\{x,y\}$.  Their product is safe.  If either rectangle is safe,
it was kept and still witnesses the edge at the new disjunction.
Otherwise both rectangles are unsafe.
Safety on assignments with sole $X$-variable $x$ implies
\[
 B_i\subseteq N(x)\quad\text{or}\quad B_j\subseteq N(x):
\]
if neither inclusion held, choosing one nonneighbor from each
$Y$-support would make their product true on an independent set.
In the first case the unique selected centre of $R_i$ is $x$;
in the second the unique selected centre of $R_j$ is $x$.
Thus the new disjunction accepts $xy$.  Every edge is preserved,
and safety proves equality with $f_G$.

Propagate the introduced zeros, bypass OR gates with a zero operand,
and remove unused gates; no new constant gate need remain.
The circuit is an OR of bilinear star products, hence
its arithmetic polynomial is quadratic with exactly the required
edge support.  Apply the exact quadratic theorem, cited in Theorem~\ref{thm:granted}.
Finally, in a $C_4$-free graph two universal rows would force
$|B_i|\leq1$, proving the stated sufficient condition.
\end{proof}

In particular, within this architecture a counterexample must contain
an unsafe rectangle $A\times\{y\}$ with at least two vertices of
$A$ adjacent to $y$, and, symmetrically, an unsafe rectangle
$\{x\}\times B$ with at least two vertices of $B$ adjacent to $x$.
Each of these rectangles also contains a nonneighbor, by unsafety.
They are necessarily distinct.
This theorem does not include mixed-degree operands or all circuits
with two levels of AND gates.

\begin{lemma}[Disjoint bad endpoint sets]
Let $P,Q$ be positive bilinear polynomials and suppose $PQ$ is safe.
The nonedge pairs occurring in $P$ and those occurring in $Q$ have
disjoint endpoint sets on each side of the bipartition.
\end{lemma}
\begin{proof}
If $xy$ occurs in $P$, $xz$ occurs in $Q$, and both pairs are
nonedges, the product contains $x^2yz$, whose support is independent.
This contradicts safety.  The argument on the other side is the same.
No four-cycle assumption is needed.
\end{proof}

\begin{example}[Both factors can be unsafe]
Let $G$ be the matching with edges $x_1y_1,x_2y_2$, and put
\[
 P=x_1y_1+x_1y_2+x_2y_2,\qquad
 Q=x_1y_1+x_2y_1+x_2y_2.
\]
Both factors contain a nonedge monomial, but their product is safe
and witnesses both matching edges.  Indeed,
\[
\begin{split}
PQ={}&x_1^2y_1^2+x_1^2y_1y_2+x_1x_2y_1^2
 +3x_1x_2y_1y_2\\
 &+x_1x_2y_2^2+x_2^2y_1y_2+x_2^2y_2^2.
\end{split}
\]
Every displayed term contains one of the matching edges.
Thus safety of a bilinear product does not imply safety of an
operand.  This is a counterexample to that proof shortcut, not to
the circuit lower bound.
\end{example}


\section{Projective completion and literal restriction}
This section counts what survives when the standard projective star
circuit is restricted to its affine subgraph, and compares that count
with a star circuit built directly for the affine graph.
\begin{proposition}[Literal restriction of the line-star circuit]
Let $q$ be a prime power and $N=q^2+q+1$.  The line-centred star
circuit for the incidence edge function of $\operatorname{PG}(2,q)$
has $N(q+2)-1$ gates.  Set the variable of one line $\ell_\infty$
and all its $q+1$ point variables to zero.  Propagating zeros in this
circuit leaves exactly
\[
 q(q+1)^2-1
\]
gates, which is $q$ more than the affine point-star upper bound
$q^2(q+2)-1$.  The literal deletion count is $q^2+2q+2$.
\end{proposition}
\begin{proof}
Each projective line has a neighborhood OR tree with $q$ gates,
followed by one AND; the outer OR tree has $N-1$ gates.  Every line
other than $\ell_\infty$ loses exactly its unique infinity point.
Its neighborhood tree therefore retains $q-1$ OR gates and its AND.
There are $N-1=q(q+1)$ surviving lines and $N-2$ outer OR gates, giving
\[
 (N-1)(q-1)+(N-1)+(N-2)=q(q+1)^2-1.
\]
Subtracting from $N(q+2)-1$ gives $q^2+2q+2$.
The affine point-star circuit has $q^2$ stars, each with $q$ neighborhood
ORs and one AND, and $q^2-1$ outer ORs.  Its size is $q^2(q+2)-1$,
smaller than the inherited circuit by $q$.
\end{proof}


\section{Why bounded-weight tests cannot give the conjectured bound}
This section proves the weight barrier and the one-point line-output construction from the earlier notes \citep[Theorem~C, Section~3.3; Theorem~D, Section~1.2]{FableNotes}. Their significance is that the circuit lower bound must use large independent sets, not merely the values on edges and small negative inputs.

Let $q$ be prime. Write $G_q$ for the incidence graph between the $q^2$ points $(a,b)\in\mathbb F_q^2$ and the $q^2$ nonvertical lines $b=sa+t$. It has $n=2q^2$ vertices and $q^3$ edges, and has no triangle or four-cycle. The graphs in this section omit the vertical parallel class; they differ from the full affine-plane graph considered earlier.

\begin{theorem}[The weight barrier]\label{thm:weightbarrier}
There is a binary monotone circuit of size $8q^2-4q-1$ agreeing with $f_{G_q}$ on every input of weight at most two. For every integer $1\le w\le2q^{1/8}$, there is such a circuit of size at most $28wn$ which accepts every input containing an edge and rejects every independent set of at most $w$ vertices.
\end{theorem}
\begin{proof}
An incidence-preserving change of coordinates is specified by
$\alpha,\gamma\in\mathbb F_q^*$ and $\beta,\delta,\epsilon\in\mathbb F_q$:
\[
 a'=\alpha a+\beta,\quad b'=\gamma b+\delta a+\epsilon,
 \quad s'=(\gamma s+\delta)/\alpha,
 \quad t'=\gamma t+\epsilon-s'\beta.
\]
Indeed $b'-s'a'-t'=\gamma(b-sa-t)$.
For one such change, compute the disjunctions $X_a,Y_b$ of the selected points with the indicated first or second coordinate, and $S_s,T_t$ of the selected lines with the indicated slope or intercept. Each of the four families uses $q(q-1)$ gates. Then compute
\[
 Z_z=\bigvee_{as=z}(X_a\wedge S_s),\qquad
 W_u=\bigvee_{z+t=u}(Z_z\wedge T_t),\qquad
 D=\bigvee_u(W_u\wedge Y_u).
\]
Every field element occurs as a product and a sum, so the first two stages each use $q^2$ products and $q^2-q$ additions. The final stage uses $q$ products and $q-1$ additions. Their total is $8q^2-4q-1$.

If the selected point set is $A$ and the selected line set is $B$, this circuit accepts precisely when some $p_1,p_2\in A$ and $\ell_1,\ell_2\in B$ satisfy
\[
 b'(p_2)=s'(\ell_1)a'(p_1)+t'(\ell_2).
\]
Every incidence satisfies this identity by repeating its point and line. If either set is empty it fails, and if both have size one it is exactly incidence. This proves the first assertion.

Choose the coordinate change uniformly. For a fixed quadruple, its acceptance identity becomes
\[
 \gamma(b_2-s_1a_1-t_2)+\delta(a_2-a_1)
       -\frac{\beta\gamma}{\alpha}(s_1-s_2)=0.
\]
If $a_1\ne a_2$, exactly one choice of $\delta$ works, conditional on the other parameters. If $a_1=a_2$ and $s_1\ne s_2$, exactly one choice of $\beta$ works. If both pairs are equal, the identity would assert that $p_2$ lies on $\ell_2$, which is impossible for an independent set. Thus, for an independent set with $|A|=a$, $|B|=b$ and $a+b\le w$, the acceptance probability is at most
\[
 \frac{a^2b^2}{q}\le\frac{w^4}{16q}\le q^{-1/2}.
\]
There are at most $n^w\le q^{3w}$ relevant nonempty sets of size at most $w$: encode a sorted set by a length-$w$ tuple, repeating its final entry. Take the conjunction of circuits from $7w$ independent coordinate choices. Its expected number of accepted independent sets of size at most $w$ is at most $q^{3w-7w/2}<1$. Some choice therefore accepts none. It still accepts every edge-containing input, and uses at most
\[
 7w(8q^2-4q-1)+(7w-1)<56wq^2=28wn
\]
gates.
\end{proof}

The result is an existence statement for nonuniform circuits. It does not give an efficient algorithm to find the successful coordinate changes. For each fixed $w$, testing correctness only on edges and independent sets of at most $w$ vertices cannot prove a lower bound proportional to $q^3$ as $q$ grows.

\begin{proposition}[All affine line outputs on one-point inputs]\label{prop:lightlines}
For every prime $q$, all $q^2+q$ affine-plane line disjunctions can be computed correctly on inputs containing at most one selected point by one binary monotone circuit of size $O(q^2)$.
\end{proposition}
\begin{proof}
Compute the $2q$ coordinate disjunctions $X_a,Y_b$ at cost $2q(q-1)$. Use the usual binary representatives $0,\ldots,q-1$. For each bit position and each bit value, take the disjunction of the corresponding coordinate wires. There are $O(\log q)$ such disjunctions, of total size $O(q\log q)$. On a one-point input they give both a bit and its complement for each coordinate; on the zero input all these wires are zero.

For each fixed slope $s$, binary multiplication by $s$, subtraction and reduction modulo $q$ compute $t=b-sa$ using $O((\log q)^2)$ Boolean gates, by the ordinary shift-and-add and long-division algorithms. Replace each wire by a pair carrying its value and complement. An AND becomes AND on the first rail and OR on the second; an OR becomes OR on the first and AND on the second; a NOT swaps the rails. Propagate all constants in the resulting circuit. The output bits of $t$ and their complements are nonconstant functions of $(a,b)$, so this produces only variable-input AND and OR gates for these outputs.

A binary prefix tree decodes the $q$ valid values of $t$. Each internal prefix extends its predecessor by one AND with a bit rail. The full tree has fewer than $2^{\lceil\log_2q\rceil+1}\le4q$ nodes, so the decoder costs $O(q)$ gates. Its leaf for $t_0$ is exactly the disjunction for $b=sa+t_0$ on a one-point input. Every gate in the constant-propagated positive circuit is zero when all its rail inputs are zero, so the zero input is also correct. The vertical outputs are already the $X_a$. Summing over the $q$ slopes gives $O(q^2+q(\log q)^2)=O(q^2)$ gates.
\end{proof}

\section{Exact searches on incidence graphs}
The searches in this section test concrete graphs below their star-circuit sizes. Every reported search stopped without a satisfying circuit or an unsatisfiability proof, so none supplies an additional lower bound.

For $\operatorname{PG}(2,3)$ there are thirteen points and thirteen lines, each of degree four. Every cover therefore has at least $52/4=13$ vertices, and either bipartition class attains thirteen. The star count is sixty-four. For the full affine plane of order three, the nine points have degree four and the twelve lines degree three; the same counting argument gives cover number nine and star count forty-four.

The encoder chooses a topological list of binary gates, their AND/OR types and their two earlier input wires. Correctness is imposed on every edge assignment and every maximal independent set. This is sufficient: every positive input contains an edge, and every negative input extends to a maximal independent set, so monotonicity determines all other values. Input permutations by graph automorphisms reduce the first gate's two inputs to orbit representatives; the rest of the circuit is unrestricted by symmetry.

The one-sided theorem removes circuits with marked degree one, and the three-missing-witness theorem removes circuits having fewer than four missing quadratic monomials. Formal degree-one and degree-two occurrence bits track these conditions without expanding the full polynomial. The terminal-product and common- or private-mask exclusions of \citet[Sections~9--10]{OneSidedCompanion} are applied to recognized circuit configurations as additional pruning. There was no candidate requiring these later lazy exclusions in the reported attempts.

\begin{center}
\small
\begin{tabular}{@{}lrrrr@{}}
\toprule
Graph & Gates tested & Negative tests & Symmetries used & Result\\\midrule
Full affine plane of order $3$ &43&197&432&undecided\\
Projective plane of order $3$ &63&704&78&undecided\\
Subdivided Petersen graph &38&579&120&undecided\\
Tutte--Coxeter graph &58&3264&720&undecided\\
\bottomrule
\end{tabular}
\end{center}
Each attempt allowed three hundred seconds of solver time after encoding. The corresponding CNF sizes, in variables and clauses, were
\[
\begin{array}{c|r@{\ ,\ }r}
\text{affine}&54\,910&1\,722\,686\\
\text{projective}&193\,413&9\,748\,983\\
\text{subdivided Petersen}&90\,164&3\,347\,138\\
\text{Tutte--Coxeter}&619\,894&35\,501\,275.
\end{array}
\]
The two latter star counts are thirty-nine and fifty-nine. The original Petersen graph has two disjoint five-cycles covering its ten vertices. Assign to each vertex the next edge on its cycle; the ten distinct assigned edges become distinct neighbors after subdivision. This gives a matching saturating the original vertices, so the subdivided graph has cover number ten and its thirty edges give the star count thirty-nine. The cubic bipartite Tutte--Coxeter graph has fifteen vertices on each side; the regular bipartite incidence count gives cover number fifteen and star count fifty-nine.

As a calibration, fixing the explicit forty-four-gate affine star circuit made the encoding satisfiable, and its decoded gate list was independently evaluated on every critical assignment. This checks a positive example; it does not certify the search encodings as infeasibility proofs. The useful-gate convention tests exactly the displayed gate count, so even an eventual unsatisfiability result would need a separate justification before it covered every smaller count. The program \emph{Incidence-circuit synthesis}, its commands, encodings and results accompany the workbench.

\section{A threshold for one negation}
Two monotone restrictions give a lower bound for circuits with one NOT gate.

Write \mbox{$(a-1)_+=\max\{a-1,0\}$} for a nonnegative integer $a$.
Let $S_1,\ldots,S_m$ be subsets of a finite input set $V$, with distinct sets meeting in at most one point. Their designated outputs are the disjunctions $\bigvee_{v\in S_i}x_v$.
The following theorem is conditional on the stated monotone lower bound; this separates the restriction argument from its base case.

\begin{theorem}\label{thm:one-not-threshold}
Assume that every family of subsets of any finite set with pairwise intersection at most one requires at least
$\sum_i(|S_i|-1)_+$ gates to compute all its disjunctions by a monotone circuit of fan-in at most two.
Let $C$ be an AND/OR/NOT circuit of fan-in at most two computing the disjunctions of a family $S_1,\ldots,S_m\subseteq V$ with this intersection property. Suppose $C$ has at most one NOT gate, and put $n=|V|$.
There is an integer $W$, $0\leq W\leq n+1$, such that
\begin{enumerate}
\item if $W\geq1$, some $a\subseteq V$ of size $W-1$ satisfies
\[
 \sum_i(|S_i\cap a|-1)_+\leq |C|;
\]
\item for every $b\subseteq V$ of size $W$,
\[
 \sum_{i:S_i\cap b=\varnothing}(|S_i|-1)_+\leq |C|.
\]
\end{enumerate}
For $W=n+1$ the second condition is vacuous.
\end{theorem}
\begin{proof}
If there is no NOT gate, take $W=0$ and apply the assumed bound to the whole family. Otherwise let $h(x)$ be the Boolean value on the input wire of the unique NOT gate. It is monotone, because its preceding circuit has no negation.
If $h$ is identically one, put $W=0$. Otherwise put
\[
 W=1+\max\{|a|:a\subseteq V,\ h(\mathbf1_a)=0\}.
\]
Choose a maximizing $a$. Monotonicity gives $h(x)=0$ on the down-set $x\leq\mathbf1_a$. The definition gives $h(\mathbf1_b)=1$ for every set $b$ of size $W$, and then $h(x)=1$ on the up-set $x\geq\mathbf1_b$. If $h$ is identically zero, this construction gives $W=n+1$.

Fix all inputs outside $a$ to zero. The NOT output becomes constant one, and the required outputs become the disjunctions of $S_i\cap a$. Propagate constants in topological order: a gate with a forcing constant becomes a constant; a gate with an identity constant redirects to its other operand; a gate with two live operands is kept. After removing designated outputs of size zero or one, every remaining output is nonconstant. Thus no constant-output gate needs to be introduced, and the resulting monotone circuit has at most $|C|$ gates. The base hypothesis applies, since intersections cannot increase. Discarded sets contribute zero to the asserted sum. This proves the first inequality.

For a fixed $b$ of size $W$, set its inputs to one and keep only the outputs with $S_i\cap b=\varnothing$. The NOT output is constant zero. The same propagation gives a monotone circuit of no greater size for the disjunctions of these sets on $V\setminus b$. Again discard sets of size at most one and apply the base hypothesis. This proves the second inequality, including $W=0$ when $h$ is identically one.
\end{proof}

\begin{corollary}\label{cor:threshold-bound}
Under the theorem's hypothesis, if a number $K$ has the following property for every $W\in\{0,\ldots,n+1\}$, then $|C|\geq K$: either $W\geq1$ and every $(W-1)$-element set $a$ has first sum at least $K$, or some $W$-element set $b$ has second sum at least $K$.
\end{corollary}
\begin{proof}
Take the value $W$ from Theorem~\ref{thm:one-not-threshold}. Its witnessing set $a$ gives the conclusion in the first alternative; its inequality for every $b$ gives the conclusion in the second.
\end{proof}

For a projective plane of order $q$ there are $N=q^2+q+1$ lines, each of size $q+1$. Let $e(A)$ be the number of lines disjoint from a point set $A$. Incidence counting gives
\[
 A_k:=\min_{|a|=k}\sum_{\ell}(|\ell\cap a|-1)_+
       =(q+1)k-N+\min_{|a|=k}e(a),
 \qquad B_k:=q\max_{|b|=k}e(b).
\]
In the threshold expression take $A_{-1}=B_{N+1}=0$.

\begin{proposition}\label{prop:plane-four-threshold}
For the Desarguesian projective plane of order four,
\[
 \min_{0\leq W\leq22}\max\{A_{W-1},B_W\}=14.
\]
Thus the preceding restriction principle yields a fourteen-gate lower bound for its line disjunctions with one NOT gate, conditional on its monotone base hypothesis. The number fourteen is the exact value of this threshold expression; it is not asserted to be the optimum circuit size or the strongest available lower bound.
\end{proposition}
\begin{proof}
There are twenty-one points and lines, with five incidences at each point. Four lines with no three concurrent have union of size $4\cdot5-\binom42=14$. Their seven-point complement avoids all four lines. Every subset of this complement also does so; hence $B_W\geq16$ for $W\leq7$.
For $W\geq8$, incidence counting and $e(a)\geq0$ give
$A_{W-1}\geq5(W-1)-21\geq14$. These two observations prove the lower bound fourteen for the minimum.

At $W=8$, the seven points of the embedded plane $\mathbb P^2(\mathbb F_2)\subseteq\mathbb P^2(\mathbb F_4)$ meet every line. Indeed a linear form over $\mathbb F_4$ restricts to an $\mathbb F_2$-linear map $\mathbb F_2^3\to\mathbb F_4$, whose kernel has a nonzero vector. Thus $A_7=35-21=14$.
Any four distinct lines have union of size at least fourteen: when adding them one at a time, the new line loses at most its number of predecessors in intersections. If an eight-point set avoided four lines, their union would lie in its thirteen-point complement, a contradiction. Therefore $B_8\leq12$. Three nonconcurrent lines have union of size twelve, so an eight-point subset of their nine-point complement shows $B_8=12$. Consequently the value at $W=8$ is fourteen, proving equality.
\end{proof}

A separate exhaustive calculation of the finite-plane subset calculation gives threshold values $2$, $6$ and $14$ for orders $2$, $3$ and $4$, respectively. At order four it checks all $2^{21}$ point subsets and finds $168$ six-point hyperovals. The geometric proof above establishes the value fourteen without relying on this enumeration.

\section{Singer convolution and the negation budget}
This section computes all projective-plane line disjunctions by cyclic
convolution and accounts for the negations needed to share that computation.
Let $N=q^2+q+1$, and let $C_s(q)$ be the least size of a fan-in-two
AND/OR/NOT circuit computing all line ORs of $\operatorname{PG}(2,q)$
with at most $s$ NOT gates. Outputs and fan-out are free.

\subsection{Singer coordinates}
The cyclic incidence description has the following direct finite-field
proof; historically it is Singer's construction \cite{Singer1938}.
Put $K=\mathbb F_{q^3}$ and identify projective points with
$K^\times/\mathbb F_q^\times$, a cyclic group of order $N$.
Choose a generator $\alpha$ of $K^\times$.
The lines are the projectivized kernels of
$x\mapsto\operatorname{Tr}_{K/\mathbb F_q}(\beta x)$, for
$\beta\in K^\times$. Indeed, trace is a nonzero linear functional:
the nonzero polynomial $T+T^q+T^{q^2}$ has degree less than $q^3$,
so it cannot vanish everywhere on $K$. Multiplication by nonzero
$\beta$ is invertible; hence the trace pairing is nondegenerate and
exhausts the three-dimensional dual space.
Two nonzero functionals have the same kernel exactly when their
ratio is a scalar in $\mathbb F_q^\times$.

Let $D=\{d\in\mathbb Z/N:\operatorname{Tr}(\alpha^d)=0\}$.
The trace kernel has $q^2-1$ nonzero vectors, so $|D|=q+1$.
Multiplication by $\alpha$ acts regularly on points and lines,
and the line sets are precisely $D+w$, $w\in\mathbb Z/N$.
Distinct lines meet once; in particular
$|D\cap(D+w)|=1$ for $w\ne0$.

\subsection{Convolution and a shared input inverter}

In Singer coordinates the outputs are
$y_w=\bigvee_{d\in D}x_{w+d}$ on $\mathbb Z/N$, with $|D|=q+1$.
For $L=\lceil\log_2(q+2)\rceil$ and $R=2^L$, form
\[
 X=\sum_{i=0}^{N-1}x_iR^i,\qquad B=\sum_{d\in D}R^{N-d}.
\]
Every base-$R$ coefficient of $XB$ is at most $q+1<R$, and the
coefficients at positions $w$ and $w+N$ sum to the number of true
inputs on line $w$. One integer multiplication followed by short
additions and nonzero tests gives
\[
 \mathcal S_N=O(N\log^2N\log\log N)
\]
gates. The network form of the multiplication bound is explicitly
part of Sch\"onhage--Strassen \cite[Summary]{SchonhageStrassen1971}.
The factors have at most $(N+1)L$ bits; their logical-network
multiplication bound is $O(t\log t\log\log t)$ for $t$-bit integers.
The $N$ block additions and nonzero tests cost $O(NL)$ further
gates. An arbitrary binary Boolean gate can be replaced by a
constant number of gates in the AND/OR/NOT basis.

Beals--Nishino--Tanaka \cite[Theorem~1.1, p.~1335;
Corollary~1.2, p.~1336]{BealsNishinoTanaka1998}
give an $m$-input inverter of size $O(m\log m)$ with
$b(m)=\lceil\log_2(m+1)\rceil$ negations, and the simulation
$C^{b(m)}(F)\leq2C(F)+O(m\log m)$ for arbitrary systems $F$.
The dual-rail implementation is explicit. Supply each input
$x_i$ and $\neg x_i$ using one shared $m$-input inverter.
If $a^+,a^-$ and $b^+,b^-$ are complementary pairs, implement an AND
by $(a^+\wedge b^+,a^-\vee b^-)$ and an OR by
$(a^+\vee b^+,a^-\wedge b^-)$.
A NOT merely interchanges two references.
Every original AND or OR costs at most two gates, original NOT
gates cost none, and positive wires at all outputs compute the
whole system. Thus the negation count is solely that of the
shared input inverter, independently of the number of outputs.
Explicitly counted constant gates contribute at most a fixed
additive allowance.

Fischer's older simulation has the same structure but a larger
inverter bound. His Theorem~16, printed p.~8 of the MIT memorandum
\cite{Fischer1975}, gives
\[
 NC_{b(m)}(F)\le 2C^\Delta(F)+O(m^2\log^2m)
              \le 6C(F)+O(m^2\log^2m),
\]
where $C^\Delta$ uses AND/OR/NOT and the second $C$ uses the full
unary/binary basis. The sharper additive term used here belongs
to Beals--Nishino--Tanaka.

\subsection{An endpoint refinement}


\begin{proposition}
For every prime power $q\geq2$,
\[
 C_{\lceil\log_2(q^2)\rceil}(q)
 \leq O(N\log^2N\log\log N).
\]
In particular, if $q=2^r$, the same bound holds with exactly the
budget $s=2r=\lfloor\log_2N\rfloor$.
\end{proposition}
\begin{proof}
Leave unprocessed a set $R_0$ consisting of one line and one point
outside it. Its complement has $m=q^2-1$ points. Use the fast
convolution circuit with all inputs in $R_0$ fixed to zero, and
apply the inverter simulation only to the $m$ remaining inputs.
This costs $2\mathcal S_N+O(m\log m)$ gates and
$b(m)=\lceil\log_2(q^2)\rceil$ negations.

The contributions from $R_0$ require $q$ OR gates on its complete
line, and one OR gate on each of the $q+1$ lines through the extra
point. All other contributions are single input wires.
Joining these with the fast contributions uses $N-1$ further
OR gates. The added cost is therefore $N+2q$, and no new negation
is used.
\end{proof}
More generally, choosing any $1\leq m\leq\min(N,2^s-1)$ inputs for the fast
part gives
\[
 C_s(q)\leq
 \min\{Nq,\ 2\mathcal S_N+O(m\log(m+1))
 +(N-m)(q+1)-\lambda(R_0)\},
\]
where $\lambda(R_0)$ counts the complete lines in the complementary
point set. A constant additive allowance handles explicit constant
gates. This is an upper envelope, not a determination of $C_s$.

If a fixed $c>0$ satisfied $C_s(q)\geq c^s qN$ for all prime-power
$q$ and all integer $s\leq\log_2N$, the proposition applied at
$q=2^r$, $s=2r$ would imply
$(2c^2)^r=O(r^2\log r)$. Hence $c\leq2^{-1/2}$.
This argument respects the integer budget; the generic simulation
at $\lceil\log_2(N+1)\rceil$ alone does not justify comparison with a
lower bound assumed only through $\lfloor\log_2N\rfloor$.

\section{Open cases}
The unrestricted conjecture remains open here. For exactly four missing quadratic witnesses, Theorem~\ref{thm:fourmaxgain} leaves only cover gain two with five deletions and an optimal quadratic frame, or cover gain three with $(D,e)=(5,0),(5,1),(6,0)$. Maximal cover gain four is settled, and every remaining gain-three case has at least two deleted additions. In the row $(d,D,e)=(3,5,0)$, Proposition~\ref{prop:pathfive} also excludes a missing $P_4\sqcup K_2$ or $K_3\sqcup K_2$. For girth at least five, Proposition~\ref{prop:pathsix} settles the missing $P_4\sqcup K_2$ shape for every optimal quadratic frame at gain three; only its excess-one case remains among the displayed alternatives. Other missing-edge shapes still require new arguments because the sixteen-transversal matching obstruction and the disjoint-context cover argument do not apply directly. General syntactic degree at most two in each side of a bipartition is also unresolved. The incidence graph of the projective plane of order three has $52$ edges and a minimum cover of size $13$, so the star construction uses $64$ gates; a circuit with $63$ gates would refute the proposed exact count for that graph. A search restricted to a symmetry class or a bounded set of monomials cannot rule out an unrestricted circuit.

\section*{Method and records}
The mathematics and text were produced by GPT-6 Astra Pro. Every accepted lower-bound argument and the scope of its separate checks are recorded at \url{\recordurl}; no human has checked the paper. Programs and certificates are available through the workbench at \url{\workbenchurl}. This text is dedicated to the public domain under CC0~1.0.

\begin{thebibliography}{99}
\bibitem[Beals et~al.(1998)]{BealsNishinoTanaka1998} Beals, R., Nishino, T., \& Tanaka, K. (1998). On the complexity of negation-limited Boolean networks. \emph{SIAM Journal on Computing, 27}(5), 1334--1347. \url{https://doi.org/10.1137/S0097539794275136}.
\bibitem[Fischer(1975)]{Fischer1975} Fischer, M. J. (1975). The complexity of negation-limited networks: A brief survey. In \emph{Automata theory and formal languages} (LNCS 33, pp. 71--82). Springer. \url{https://doi.org/10.1007/3-540-07407-4_9}. Locators refer to MIT memorandum TM-065, Theorem~16 and Lemmas~18--19, pp. 8--10.
\bibitem[Fable notes(2026)]{FableNotes} \emph{The hard links of the P versus NP programme: Connections, partial results and obstructions}; \emph{Implications for the proof graph}. (2026). Unpublished mathematical notes, 23 November 2026. Theorem~C, Section~3.3 of the first note; Theorem~D, Section~1.2 of the second. Proofs of these two statements are given in this paper.
\bibitem[Jerrum and Snir(1982)]{JerrumSnir1982} Jerrum, M., \& Snir, M. (1982). Some exact complexity results for straight-line computations over semirings. \emph{Journal of the ACM, 29}(3), 874--897. \url{https://doi.org/10.1145/322326.322341}.
\bibitem[Jukna(2005)]{JuknaSingle} Jukna, S. (2005). \emph{Disproving the single level conjecture}. Electronic Colloquium on Computational Complexity, TR05-021, revision~2. \url{https://eccc.weizmann.ac.il/report/2005/021/}.
\bibitem[Jukna(2025)]{JuknaRead} Jukna, S. (2025). Notes on Boolean read-$k$ and multilinear circuits. \emph{Discrete Applied Mathematics, 360}, 307--327. \url{https://doi.org/10.1016/j.dam.2024.09.023}. Locators here refer to the arXiv manuscript \url{https://arxiv.org/abs/2207.08701}.
\bibitem[One-sided companion paper(2026)]{OneSidedCompanion} \emph{The one-sided edge-circuit bound of pequalnp.com and circuit lifts to inverse semirings, after Faul, Goswami, Joubert and Manuell}. (2026). Unpublished companion manuscript, 24 November 2026, Theorem~1.1 and Sections~9--10. Records at \url{\recordurl}.
\bibitem[Companion paper(2026)]{Companion} \emph{The P = NP? edge-circuit conjecture at degree one, and zero-adjoined rings after Faul, Goswami, Joubert and Manuell}. (2026). Unpublished companion manuscript, 24 November 2026. Records at \url{\recordurl}.
\bibitem[P versus NP programme(2026)]{Fleet} P versus NP programme. (2026). \emph{Public proof graph and Lean repository}. \url{https://pequalnp.com}; \url{https://github.com/brilliant-hq/pequalnp-proofs}.
\bibitem[Sch\"onhage and Strassen(1971)]{SchonhageStrassen1971} Sch\"onhage, A., \& Strassen, V. (1971). Schnelle Multiplikation gro{\ss}er Zahlen. \emph{Computing, 7}, 281--292. \url{https://doi.org/10.1007/BF02242355}.
\bibitem[Singer(1938)]{Singer1938} Singer, J. (1938). A theorem in finite projective geometry and some applications to number theory. \emph{Transactions of the American Mathematical Society, 43}(3), 377--385. \url{https://doi.org/10.1090/S0002-9947-1938-1501951-4}.
\end{thebibliography}
\end{document}
